\chapter{Endocrinology}

\medterm{Acromegaly}-- Chronic disorder of excess growth hormone secretion after epiphyseal closure in adults causing progressive enlargement of acral structures including hands, feet, jaw, and soft tissues. Most commonly caused by a pituitary somatotroph adenoma accounting for approximately ninety-five percent of cases. Clinical features include coarsened facial features, prognathism, macroglossia, widening of the hands and feet, soft tissue swelling causing carpal tunnel syndrome, excessive sweating, hypertension, diabetes mellitus, and cardiovascular disease. Diagnosis is confirmed by elevated IGF-1 levels and failure of growth hormone suppression during oral glucose tolerance testing. Treatment options include transsphenoidal surgery for adenoma resection, medical therapy with somatostatin analogues, dopamine agonists, or growth hormone receptor antagonists, and radiation therapy for residual or refractory disease.

\medterm{Acromegaly Complications}-- Systemic complications of chronic growth hormone excess include diabetes mellitus from insulin resistance, cardiovascular disease with cardiomyopathy and arrhythmias, sleep apnea from macroglossia and upper airway soft tissue hypertrophy, hypertension, colon polyps and increased colorectal cancer risk, osteoarthritis, carpal tunnel syndrome, and increased mortality. Screening colonoscopy is recommended at diagnosis and at regular intervals. Cardiac assessment with echocardiography is important. Regular dermatological surveillance for skin lesions and colonoscopic screening according to standard guidelines are recommended.

\medterm{Addison Disease}-- Primary adrenal insufficiency resulting from destruction or dysfunction of the adrenal cortex, leading to deficiency of cortisol, aldosterone, and adrenal androgens. Most common cause in developed countries is autoimmune adrenalitis (autoimmune polyglandular syndrome type 2), with anti-21-hydroxylase antibodies present in 85-90 percent of cases. Other causes include infections (tuberculosis, histoplasmosis, CMV in HIV/AIDS), hemorrhage, metastatic infiltration, and genetic disorders (congenital adrenal hyperplasia, adrenoleukodystrophy, and autoimmune polyglandular syndrome type 1). Presents with fatigue, hyperpigmentation of the skin and oral mucosa, weight loss, hypotension, and salt craving. Diagnosis is made by elevated ACTH with a low cortisol level and a poor response to cosyntropin stimulation testing. Treatment requires lifelong glucocorticoid and mineralocorticoid replacement therapy with stress dose adjustments during illness or surgery.

\medterm{Addisonian Crisis}-- Life-threatening emergency resulting from acute cortisol deficiency in patients with adrenal insufficiency. Precipitated by infection, trauma, surgery, or abrupt discontinuation of chronic glucocorticoid therapy. Presents with severe hypotension or shock unresponsive to fluids and vasopressors, profound weakness, nausea, vomiting, abdominal pain, fever, confusion, and hypoglycemia. Laboratory findings include hyponatremia, hyperkalemia, hypoglycemia, and eosinophilia. Treatment is immediate IV hydrocortisone, aggressive fluid resuscitation with normal saline and dextrose to correct hypoglycaemia, vasopressor support if needed, and treatment of the underlying precipitating cause.

\medterm{Adrenal Cortex Physiology}-- The adrenal cortex produces three classes of steroid hormones from the three zones. The zona glomerulosa produces mineralocorticoids primarily aldosterone regulating sodium and potassium balance under the influence of angiotensin II and potassium. The zona fasciculata produces glucocorticoids primarily cortisol regulating glucose metabolism protein catabolism and the stress response under the influence of ACTH. The zona reticularis produces adrenal androgens including DHEA and DHEA-S under the influence of ACTH.

\medterm{Adrenal Hemorrhage}-- Bleeding into the adrenal gland most commonly occurring in critically ill patients with sepsis, coagulopathy, or anticoagulation therapy. Bilateral adrenal hemorrhage may cause acute adrenal insufficiency with circulatory collapse. Waterhouse-Friderichsen syndrome: adrenal hemorrhage associated with meningococcal sepsis. Diagnosis is by CT showing enlarged adrenal glands with hemorrhage. Treatment includes stress-dose corticosteroids for bilateral involvement.

\medterm{Adrenal Insufficiency}-- Deficient production of cortisol by the adrenal cortex classified as primary adrenal insufficiency or Addison disease from adrenal gland destruction, secondary from pituitary ACTH deficiency, and tertiary from chronic exogenous glucocorticoid suppression of the hypothalamic-pituitary-adrenal axis. Primary causes include autoimmune adrenalitis accounting for eighty percent in developed countries, tuberculosis, fungal infections, metastatic disease, and hemorrhage. Presents with fatigue, weakness, weight loss, hyperpigmentation, hypotension, and salt craving in primary disease, whereas secondary deficiency presents without hyperpigmentation or hyperkalaemia.

\medterm{Adrenal Medulla Physiology}-- The adrenal medulla produces catecholamines primarily epinephrine and norepinephrine in response to sympathetic nervous system activation. Epinephrine accounts for approximately eighty percent and norepinephrine approximately twenty percent of catecholamine output. These hormones mediate the fight-or-flight response through alpha and beta adrenergic receptors causing tachycardia hypertension bronchodilation and metabolic mobilization including glycogenolysis and lipolysis. Pheochromocytomas arise from the adrenal medulla and secrete excess catecholamines causing episodic hypertension, headache, sweating, and palpitations.

\medterm{Adrenal Myelolipoma}-- Benign adrenal tumor composed of mature adipose tissue and hematopoietic elements. Usually asymptomatic and discovered incidentally. No malignant potential. May cause adrenal hemorrhage if large. Diagnosis is by imaging showing fat attenuation on CT. No treatment required unless symptomatic from mass effect or hemorrhage.

\medterm{Adrenal Venous Sampling}-- Interventional radiology procedure for lateralizing aldosterone secretion in primary hyperaldosteronism. Catheterization of both adrenal veins allows measurement of aldosterone and cortisol levels to calculate selectivity index confirming successful cannulation and lateralization ratio. Lateralization to one gland suggests aldosterone-producing adenoma amenable to adrenalectomy. Bilateral secretion suggests bilateral adrenal hyperplasia managed medically. The procedure is technically challenging and requires an experienced interventional radiologist, with success rates of approximately seventy to eighty percent for right adrenal vein cannulation and higher for the left.

\medterm{Aicardi-Goutieres Syndrome}-- Rare autosomal recessive inflammatory disorder mimicking congenital viral infection. Characterized by cerebral calcifications, white matter disease, cerebral atrophy, and cerebrospinal fluid lymphocytosis. Presents with severe neurodevelopmental impairment, seizures, and spasticity. Some forms involve interferon pathway dysregulation with elevated interferon-alpha in cerebrospinal fluid. Diagnosis is by clinical features, brain imaging showing calcifications, and genetic testing. No specific treatment is available; management is supportive and includes anticonvulsants, physiotherapy, and management of feeding difficulties. Prognosis is poor with many children not surviving beyond early childhood.

\medterm{Alpha-Glucosidase Inhibitors}-- Oral glucose-lowering medications including acarbose, miglitol, and voglibose that delay carbohydrate digestion and glucose absorption in the small intestine by inhibiting brush border alpha-glucoside enzymes. Reduce postprandial glucose excursions without affecting fasting glucose significantly. Side effects include flatulence, bloating, and diarrhea from undigested carbohydrates reaching the colon. Contraindicated in inflammatory bowel disease and hepatorenal impairment. Effective A1c reduction of approximately 0.5 to 0.8 percent is modest compared with other oral agents.

\medterm{Amiodarone-Induced Thyroid Dysfunction}-- Thyroid abnormalities occurring in approximately fifteen to twenty percent of patients on amiodarone due to the high iodine content of the molecule. Type one results from iodine-induced increased thyroid hormone synthesis in patients with underlying thyroid disease. Type two results from destructive thyroiditis. Diagnosis is by thyroid function tests, thyroid uptake, and color flow Doppler. Type one is treated with thionamides. Type two is treated with corticosteroids. Amiodarone-induced hypothyroidism is managed with levothyroxine replacement without discontinuing amiodarone if it is clinically indicated. Amiodarone should be discontinued if possible when hyperthyroidism develops.

\medterm{Amylin Analogues}-- Pramlintide is a synthetic analogue of amylin a co-secreted hormone from pancreatic beta cells that slows gastric emptying suppresses postprandial glucagon secretion and promotes satiety. Used as adjunctive therapy with insulin in type one and type two diabetes. Reduces postprandial glucose excursions and A1c. Side effects include nausea and injection site reactions. Risk of hypoglycemia when combined with insulin requiring dose reduction.

\medterm{Antidiuretic Hormone Physiology}-- Antidiuretic hormone also known as vasopressin is synthesized in the supraoptic and paraventricular nuclei of the hypothalamus and stored in the posterior pituitary. Released in response to increased serum osmolality detected by hypothalamic osmoreceptors and decreased blood volume detected by baroreceptors. Acts on the collecting duct through V2 receptors inserting aquaporin-2 water channels causing water reabsorption and urine concentration. Also causes vasoconstriction through V1 receptors at higher concentrations causing vasoconstriction. Deficiency leads to diabetes insipidus, while excess causes syndrome of inappropriate antidiuretic hormone secretion.

\medterm{Aromatase Excess Syndrome}-- Rare autosomal dominant condition caused by mutations in the CYP19A1 gene leading to overexpression of aromatase enzyme and excessive estrogen production. Presents in males with gynecomastia, advanced bone age, and accelerated linear growth with early epiphyseal closure. In females with precocious puberty, macroclitoris, and virilization. Diagnosis is by elevated estradiol and suppressed gonadotropins. Treatment is with aromatase inhibitors anastrozole or letrozole to reduce estrogen synthesis and manage symptoms.

\medterm{Blood Glucose Monitoring}-- Blood sugar monitoring involves measuring glucose levels in the blood to manage diabetes, typically using a glucometer with finger-stick testing or continuous glucose monitors that track interstitial glucose in real time. Regular monitoring helps patients and clinicians adjust insulin doses, diet, and physical activity to maintain target glycemic ranges. Advances include flash glucose monitoring systems and smartphone-integrated sensors that improve convenience and data trends.

\medterm{Bone Densitometry}-- Dual-energy X-ray absorptiometry measures bone mineral density at the lumbar spine and hip. Results are expressed as T-score comparing to young adult mean and Z-score comparing to age-matched controls. T-score of minus 1.0 to minus 2.5 indicates osteopenia. T-score of minus 2.5 or less indicates osteoporosis. T-score of minus 2.5 with fragility fracture indicates severe osteoporosis. Fracture risk assessment using FRAX incorporates clinical risk factors and BMD to estimate ten-year probability of hip fracture over ten years.

\medterm{Bone Turnover Markers}-- Biochemical markers reflecting the rate of bone remodeling. Formation markers include bone-specific alkaline phosphatase, osteocalcin, and N-terminal and C-terminal propeptides of type I procollagen. Resorption markers include C-terminal and N-terminal telopeptides of type I collagen and deoxypyridinoline. Elevated markers indicate increased bone turnover as in primary hyperparathyroidism Paget disease and osteomalacia. Suppressed markers indicate low bone turnover as in adynamic bone disease. Useful for monitoring response to osteoporosis therapy and assessing metabolic bone disease.

\medterm{Calcitonin}-- A 32-amino-acid peptide hormone secreted by the parafollicular C-cells of the thyroid gland in response to elevated serum calcium. It lowers blood calcium by inhibiting osteoclast-mediated bone resorption and increasing renal calcium excretion. Its physiologic role in humans is minor compared to PTH and vitamin D; however, calcitonin is a sensitive tumor marker for medullary thyroid carcinoma (MTC). Synthetic calcitonin (salmon calcitonin, more potent than human) was historically used for Paget disease of bone, osteoporosis, and hypercalcemia but has been largely replaced by bisphosphonates, denosumab, and other agents due to modest efficacy and association with malignancy risk with long-term use.

\medterm{Calcium Sensing Receptor}-- G-protein coupled receptor expressed primarily on parathyroid glands and renal tubular cells that senses extracellular ionized calcium levels. On parathyroid glands activation suppresses PTH secretion. On renal tubular cells activation increases calcium excretion and decreases phosphate reabsorption. Inactivating mutations cause familial hypocalciuric hypercalcemia and neonatal severe hyperparathyroidism. Activating mutations cause autosomal dominant hypocalcemia. The calcium sensing receptor is the target for cinacalcet, a calcimimetic used in secondary hyperparathyroidism of chronic kidney disease and parathyroid carcinoma.

\medterm{Carcinoid Heart Disease}-- Cardiac complication of carcinoid syndrome occurring in approximately fifty percent of patients with hepatic metastases. Serotonin and other vasoactive substances bypass hepatic metabolism and cause endocardial fibrosis affecting the right-sided heart valves particularly the tricuspid and pulmonary valves. Tricuspid regurgitation is the most common finding followed by pulmonary stenosis. Presents with right heart failure including edema ascites and fatigue. Diagnosis is by echocardiography. Treatment includes diuretics for heart failure and valve replacement surgery for severe valvular disease. Somatostatin analogues may slow disease progression.

\medterm{Catecholamine Crisis}-- Acute life-threatening elevation of catecholamines causing severe hypertension, tachycardia, and multi-organ dysfunction. Occurs in pheochromocytoma and paraganglioma precipitated by anesthesia, surgery, trauma, or medications that stimulate catecholamine release. Treatment includes immediate IV phentolamine or nitroprusside for blood pressure control followed by beta-blockade once adequate alpha-blockade is established. Beta-blockers should never be given before alpha-blockade due to the risk of unopposed alpha-adrenergic stimulation causing hypertensive crisis.

\medterm{Closed-Loop Insulin Delivery}-- Automated insulin delivery systems that combine continuous glucose monitoring with an insulin pump and control algorithm. The algorithm adjusts basal insulin delivery based on sensor glucose values in real time. Hybrid closed-loop systems require user-initiated meal boluses but automatically manage basal rates. Fully automated systems are in development. Current systems include Medtronic 780G, Omnipod 5, Tandem Control-IQ, and CamAPS FX. Improve time in range and reduce hypoglycemia compared to conventional insulin pump therapy.

\medterm{Congenital Adrenal Hyperplasia}-- Group of autosomal recessive disorders caused by enzyme deficiencies in the adrenal steroidogenesis pathway. Twenty-one hydroxylase deficiency is the most common accounting for approximately ninety-five percent of cases. Impaired cortisol synthesis leads to elevated ACTH driving adrenal androgen excess causing virilization in females and precocious puberty in males. Severe forms present with salt-wasting crisis in neonates from aldosterone deficiency. Diagnosis is by elevated 17-hydroxyprogesterone levels and genetic testing. Treatment includes glucocorticoid replacement to suppress ACTH and androgen excess, mineralocorticoid replacement for salt-wasting forms, and corrective surgery for virilised genitalia in females.

\medterm{Congenital Hyperinsulinism}-- Condition of inappropriate insulin secretion in neonates and infants causing severe hypoglycemia. Most commonly caused by mutations in the KATP channel genes ABCC8 and KCNJ11. Histologically classified into focal and diffuse forms. Presents with severe symptomatic hypoglycemia in the first days of life. Diagnosis requires demonstration of inappropriately elevated insulin and C-peptide during hypoglycemia. Treatment includes continuous glucose infusion, diazoxide as first-line medical therapy, and octreotide for diazoxide-unresponsive cases. Surgical resection of focal lesions or near-total pancreatectomy for diffuse disease may be necessary.

\medterm{Congenital Hyperthyroidism}-- Thyrotoxicosis present at birth most commonly caused by maternal Graves disease with transplacental passage of thyroid-stimulating immunoglobulins. Presents with tachycardia, irritability, poor feeding, weight loss, and goiter. Transient neonatal hyperthyroidism resolves as maternal antibodies are cleared over weeks to months. Neonatal Graves disease requires treatment with methimazole and beta-blockers. Rare causes include TSH receptor mutations causing constitutive activation and McCune-Albright syndrome.

\medterm{Congenital Hypothyroidism}-- Thyroid hormone deficiency present at birth occurring in approximately one in two thousand to four thousand live births. Most common cause is thyroid dysgenesis including aplasia, ectopia, and hypoplasia. Universal newborn screening with TSH or T4 allows early detection and treatment preventing intellectual disability. Presentation in unscreened infants includes prolonged jaundice, constipation, poor feeding, lethargy, hypotonia, macroglossia, and umbilical hernia. Treatment is levothyroxine replacement starting at ten to fifteen micrograms per kilogram daily with dose adjustment based on thyroid function tests.

\medterm{Continuous Glucose Monitoring}-- Technology measuring interstitial glucose levels continuously providing real-time glucose data, trends, and alerts. Sensors are inserted subcutaneously typically in the abdomen or arm and measure glucose every five minutes. Devices include Dexcom, FreeStyle Libre, and Medtronic Guardian systems. Provides time in range the percentage of time glucose remains between seventy and one hundred eighty milligrams per deciliter, glucose variability metrics, and trend arrows. Alerts notify patients of impending or current hypoglycaemia and hyperglycaemia. Data can be shared remotely with caregivers and healthcare providers, improving diabetes management.

\medterm{Cushing Disease}-- Hypercortisolism caused by a pituitary ACTH-secreting adenoma, the most common cause of endogenous Cushing syndrome (70%). Presents with central obesity, moon face, buffalo hump, violaceous striae, proximal myopathy, easy bruising, osteoporosis, hypertension, hyperglycemia, depression, and menstrual irregularities. Diagnosis: elevated 24-hour urinary free cortisol, abnormal dexamethasone suppression test, late-night salivary cortisol; ACTH level differentiates ACTH-dependent from independent causes; pituitary MRI localizes adenoma. Treatment: transsphenoidal adenomectomy is first-line; bilateral adrenalectomy for refractory cases.

\medterm{Cushing Syndrome}-- Clinical state resulting from chronic exposure to excess glucocorticoids regardless of cause. Most common endogenous cause is Cushing disease from pituitary ACTH-secreting adenoma accounting for approximately seventy percent. Ectopic ACTH syndrome from non-pituitary tumors including small cell lung cancer and carcinoid tumors accounts for approximately fifteen percent. Adrenal adenomas and carcinomas account for approximately fifteen percent. Exogenous glucocorticoid administration is the most common cause of exogenous Cushing syndrome from prescribed corticosteroids including prednisolone and dexamethasone for inflammatory conditions.

\medterm{Cushing Syndrome Complications}-- Untreated Cushing syndrome carries significant morbidity and mortality from cardiovascular disease, infections, thromboembolic events, and metabolic derangements. Cardiovascular complications include accelerated atherosclerosis, hypertension, and dyslipidemia. Infection risk is increased from immunosuppression. Venous thromboembolism occurs at high rates. Metabolic complications include diabetes, osteoporosis, and myopathy. Psychiatric manifestations include depression, cognitive impairment, and psychosis. Cardiovascular risk remains elevated even after biochemical cure, necessitating long-term surveillance.

\medterm{Dawn Phenomenon}-- Physiological increase in blood glucose occurring in the early morning hours between approximately three and eight AM driven by counter-regulatory hormone surges including growth hormone, cortisol, and catecholamines. More pronounced in patients with diabetes who lack the compensatory insulin response. Can be identified by comparing overnight and fasting glucose values. Differentiated from the Somogyi effect which is rebound hyperglycemia after nocturnal hypoglycemia. Management involves adjustment of basal insulin rates and consideration of continuous glucose monitoring to detect and address early morning hyperglycaemia.

\medterm{Delayed Puberty}-- Failure to develop secondary sexual characteristics by age thirteen in girls and age fourteen in boys. Causes include constitutional delay which is the most common, chronic illness, malnutrition, hypogonadotropic hypogonadism from Kallmann syndrome or pituitary disease, and hypergonadotropic hypogonadism from gonadal failure. Diagnosis includes bone age assessment, hormone levels including LH, FSH, testosterone or estradiol, and imaging. Treatment is sex steroid replacement to induce puberty with monitoring of growth and bone age progression. Underlying causes should be addressed where possible.

\medterm{Denosumab for Osteoporosis}-- Fully human monoclonal antibody targeting RANK ligand preventing osteoclast formation activation and survival. Reduces fracture risk at the spine hip and non-vertebral sites. Administered as subcutaneous injection every six months. Advantages include no renal dose adjustment and effectiveness even with oral bisphosphonate intolerance. Disadvantages include risk of vertebral fractures after discontinuation due to rebound bone loss requiring bisphosphonate transition therapy or extended denosumab therapy indefinitely.

\medterm{Diabetes and Metabolic Health}-- Diabetes and metabolic health are closely linked through insulin resistance, obesity, dyslipidemia, and hypertension, collectively known as the metabolic syndrome. Poor metabolic health increases the risk of developing type 2 diabetes and worsens glycemic control in existing diabetes. Improving metabolic health through weight loss, dietary changes, and physical activity can prevent or reverse type 2 diabetes in some individuals.

\medterm{Diabetes Complications}-- Diabetes complications arise from chronic hyperglycemia damaging blood vessels and nerves over time. Microvascular complications include diabetic retinopathy leading to vision loss, nephropathy progressing to kidney failure, and neuropathy causing pain and loss of sensation. Macrovascular complications include accelerated atherosclerosis manifesting as heart attack, stroke, and peripheral artery disease.

\medterm{Diabetes insipidus}-- A disorder characterized by polyuria (excessive urine output exceeding 40 mL/kg/24 hours) and polydipsia due to deficient secretion of antidiuretic hormone (ADH, vasopressin) from the posterior pituitary (central diabetes insipidus) or renal resistance to ADH action (nephrogenic diabetes insipidus). Central DI is caused by destruction of the hypothalamic-pituitary axis from head trauma, neurosurgery, tumors (craniopharyngioma, germinoma), infiltrative diseases (Langerhans cell histiocytosis, sarcoidosis), and idiopathic causes. Nephrogenic DI is caused by genetic mutations (AVPR2 or AQP2), drugs (lithium, demeclocycline), hypercalcaemia, hypokalaemia, and chronic kidney disease. Diagnosis involves water deprivation test measuring urine osmolality and response to desmopressin. Central DI shows marked improvement with desmopressin, whereas nephrogenic DI shows minimal response. Treatment includes desmopressin for central DI, correction of underlying cause for nephrogenic DI, and thiazide diuretics with a low-sodium diet for partial nephrogenic DI.

\medterm{Diabetes Medications}-- Diabetes medications lower blood glucose through various mechanisms: metformin reduces hepatic glucose production; sulfonylureas stimulate insulin secretion; SGLT2 inhibitors increase urinary glucose excretion; GLP-1 receptor agonists enhance incretin effects; and insulin replaces deficient hormone. Choice of therapy depends on diabetes type, glycemic targets, patient preferences, cardiovascular and renal status, and cost.

\medterm{Diabetes Mellitus}-- Diabetes mellitus is a metabolic disorder characterized by chronic hyperglycemia resulting from insufficient insulin production, insulin resistance, or both. Type 1 diabetes results from autoimmune destruction of pancreatic beta cells, while type 2 is driven by genetic and lifestyle factors leading to insulin resistance. Long-term complications affect the eyes, kidneys, nerves, and cardiovascular system; management focuses on glycemic control, diet, exercise, and medications.

\medterm{Diabetes Mellitus Type 1}-- An autoimmune disease characterized by destruction of pancreatic beta cells leading to absolute insulin deficiency. Usually presents in childhood or young adulthood with acute onset of polyuria, polydipsia, weight loss, and diabetic ketoacidosis (DKA) at diagnosis in many. Laboratory: hyperglycemia, positive autoantibodies (GAD, IA-2, ZnT8), low C-peptide. Lifelong insulin therapy required. Intensive glycemic control reduces microvascular and macrovascular complications.

\medterm{Diabetes Mellitus Type 2}-- A progressive metabolic disorder characterized by insulin resistance and relative insulin deficiency, strongly associated with obesity and physical inactivity. Typically presents in adults but increasingly seen in younger populations. Diagnosis: fasting glucose \(\ge\)126 mg/dL, HbA1c \(\ge\)6.5%, or OGTT \(\ge\)200 mg/dL. Management: lifestyle modification, metformin first-line, followed by additional oral agents, GLP-1 receptor agonists, SGLT2 inhibitors, and ultimately insulin as beta-cell function declines.

\medterm{Diabetes Prevention}-- Preventing diabetes focuses on reducing the risk of developing type 2 diabetes through lifestyle interventions and risk factor management. Maintaining a healthy body weight, engaging in regular physical activity, and following a balanced diet low in refined sugars and saturated fats are cornerstone strategies. For individuals with prediabetes, structured programs involving dietary changes, exercise, and sometimes metformin can significantly lower the likelihood of progression to overt diabetes.

\medterm{Diabetic Foot Disease}-- Complex multifactorial condition resulting from the combination of peripheral neuropathy, peripheral arterial disease, and impaired immune function in diabetes. Neuropathy leads to loss of protective sensation and foot deformities. Peripheral arterial disease impairs healing. Minor trauma progresses to ulceration and infection. The Wagner classification grades ulcers from zero to five based on depth and presence of gangrene. Treatment includes offloading, wound care, infection management, revascularisation for peripheral arterial disease, and amputation for non-salvageable tissue loss. Multidisciplinary foot care teams are essential for optimal outcomes.

\medterm{Diabetic Ketoacidosis}-- An acute, life-threatening complication of diabetes mellitus, more common in type 1. Characterized by hyperglycemia (>250 mg/dL), metabolic acidosis (pH <7.3, bicarbonate <15 mEq/L), and ketonemia. Precipitating factors: infection, missed insulin doses, MI, pancreatitis, and stress. Presents with polyuria, polydipsia, nausea, vomiting, abdominal pain, Kussmaul respirations, and altered mental status. Treatment: fluid resuscitation, insulin infusion, potassium replacement, and correction of precipitating cause.

\medterm{Diabetic Ketoacidosis Monitoring}-- Monitoring during DKA management includes hourly blood glucose measurement, serum electrolytes every two to four hours, venous blood gas assessment, beta-hydroxybutyrate levels, and fluid balance. Resolution criteria include blood glucose less than two hundred milligrams per deciliter, bicarbonate greater than fifteen milliequivalents per liter, venous pH greater than seven point three, and anion gap less than twelve. Transition to subcutaneous insulin occurs when resolution criteria are met and the patient is eating and drinking normally. Frequent monitoring of potassium is essential because levels fall during insulin therapy and require aggressive replacement to prevent hypokalaemia.

\medterm{Diabetic Neuropathy}-- Most common complication of diabetes affecting approximately fifty percent of patients. Distal symmetric polyneuropathy presents with stocking-glove sensory loss, burning pain, and loss of proprioception. Autonomic neuropathy causes gastroparesis, orthostatic hypotension, erectile dysfunction, and sudomotor dysfunction. Mononeuropathies include cranial nerve palsies particularly third nerve palsy and carpal tunnel syndrome. Diagnosis is by clinical examination and nerve conduction studies. Treatment includes optimised glycaemic control, symptomatic management of neuropathic pain with agents such as gabapentin, pregabalin, amitriptyline, or duloxetine, and management of autonomic complications.

\medterm{Diabetic Retinopathy}-- A microvascular complication of diabetes mellitus and a leading cause of blindness in working-age adults. Non-proliferative DR (NPDR): microaneurysms, dot-blot hemorrhages, exudates (hard and cotton-wool spots), and venous beading. Proliferative DR (PDR): retinal neovascularization, vitreous hemorrhage, and tractional retinal detachment. Macular edema (DME) can occur at any stage and is the most common cause of vision loss. Risk factors: poor glycemic control, hypertension, duration of diabetes. Management: strict glucose and blood pressure control; anti-VEGF injections (ranibizumab, aflibercept, bevacizumab) and focal laser for DME; panretinal photocoagulation for PDR.

\medterm{Dipeptidyl Peptidase-4 Inhibitors}-- Oral glucose-lowering medications including sitagliptin, saxagliptin, linagliptin, and alogliptin that inhibit the enzyme degrading incretin hormones GLP-1 and GIP. Resulting increase in active incretin levels enhances glucose-dependent insulin secretion and suppresses glucagon. Neutral effect on weight and low hypoglycemia risk. Cardiovascular safety has been demonstrated though some agents showed increased heart failure hospitalizations. Dose adjustment required in renal impairment for all except linagliptin which is hepatobiliary excreted.

\medterm{Dwarfism}-- Short stature resulting from a medical condition, most commonly achondroplasia (the most common form of skeletal dysplasia, caused by gain-of-function mutation in FGFR3, presenting with rhizomelic limb shortening, macrocephaly, frontal bossing, and trident hands). Other causes: growth hormone deficiency, hypothyroidism, Turner syndrome, Noonan syndrome, skeletal dysplasias (osteogenesis imperfecta, spondyloepiphyseal dysplasia), and intrauterine growth restriction. Diagnosis: history, physical exam, growth chart analysis, skeletal survey, genetic testing, and endocrine evaluation. Management: growth hormone therapy for GH deficiency; limb lengthening procedures in selected cases; management of complications (spinal stenosis, obesity, obstructive sleep apnea, otitis media).

\medterm{Ectopic ACTH Syndrome}-- Adrenocorticotropic hormone secretion by non-pituitary tumors causing Cushing syndrome. Most commonly caused by small cell lung cancer and bronchial carcinoid tumors. Other sources include thymic carcinoid, pancreatic neuroendocrine tumors, and medullary thyroid carcinoma. Ectopic ACTH typically causes more severe hypercortisolism than Cushing disease with pronounced hypokalemia and metabolic alkalosis. Diagnosis requires demonstration of non-suppressible cortisol after high-dose dexamethasone suppression test, inferior petrosal sinus sampling for localisation of ACTH source, and imaging of the pituitary and adrenal glands.

\medterm{Ectopic Thyroid Tissue}-- Thyroid tissue located outside the normal pretracheal position resulting from failure of normal thyroid descent during embryological development. Most commonly found at the base of the tongue termed lingual thyroid. May be the only functioning thyroid tissue. Symptoms include dysphagia, dysphonia, and upper airway obstruction. Diagnosis is by imaging with thyroid scan confirming uptake in the ectopic location. Treatment includes levothyroxine supplementation to suppress TSH and reduce ectopic tissue size. Surgical excision is reserved for symptomatic lesions.

\medterm{Empty Sella Syndrome}-- Herniation of subarachnoid space into the sella turcica, compressing the pituitary gland and flattening it against the sellar floor. Primary: due to congenital weakness of the diaphragma sella. Secondary: post-surgery, post-radiation, or pituitary apoplexy (Sheehan syndrome). Most are asymptomatic. Symptoms when present: headache, visual field defects, and endocrine dysfunction (hypopituitarism, growth hormone deficiency, hyperprolactinemia). Diagnosis: MRI shows CSF filling the sella with compressed pituitary. Treatment: endocrine replacement as needed; no treatment for the empty sella itself.

\medterm{Endocrine Disrupting Chemicals}-- Environmental chemicals that interfere with endocrine function including bisphenol A, phthalates, pesticides, and polychlorinated biphenyls. Potential associations with diabetes, obesity, thyroid dysfunction, reproductive disorders, and hormone-sensitive cancers. Evidence is strongest for associations with type two diabetes and obesity. Avoidance strategies include limiting plastic food containers choosing organic produce and avoiding personal care products with parabens and phthalates. Regulatory bodies continue to evaluate the health impact of these chemicals and develop guidelines for safe exposure limits.

\medterm{Endocrine Glands}-- Ductless glands that secrete hormones directly into the bloodstream to regulate distant target organs. Major endocrine glands: hypothalamus (releasing/inhibiting hormones regulating the pituitary), pituitary gland (anterior: GH, TSH, ACTH, FSH, LH, prolactin; posterior: ADH, oxytocin), thyroid gland (T4, T3, calcitonin), parathyroid glands (PTH), adrenal glands (cortex: cortisol, aldosterone, androgens; medulla: epinephrine, norepinephrine), pancreas (islets of Langerhans: insulin, glucagon, somatostatin), ovaries (estrogen, progesterone), testes (testosterone), pineal gland (melatonin), and kidneys (erythropoietin, renin). The endocrine system maintains homeostasis through feedback loops (primarily negative feedback), with the hypothalamic-pituitary axis serving as the master regulator.

\medterm{Endocrinology}-- The branch of medicine concerned with the study, diagnosis, and treatment of disorders of the endocrine system--the network of glands that secrete hormones regulating metabolism, growth and development, tissue function, sexual function, reproduction, sleep, and mood. Endocrinologists manage conditions including diabetes mellitus, thyroid disorders (hyperthyroidism, hypothyroidism, nodules, cancer), adrenal disorders (Cushing syndrome, Addison disease, pheochromocytoma), pituitary disorders (acromegaly, prolactinoma, hypopituitarism), parathyroid and calcium disorders (hyperparathyroidism, osteoporosis), and reproductive endocrinology (infertility, PCOS, menopause). Diagnostic tools: hormone assays, stimulation/suppression tests, and imaging (ultrasound, CT, MRI, nuclear medicine).

\medterm{Euthyroid Sick Syndrome}-- Abnormal thyroid function tests in patients with non-thyroidal illness (critical illness, starvation, major surgery) without intrinsic thyroid disease. Pattern: low T3 (most common), low T4, reverse T3 elevated, normal or low TSH. The severity correlates with illness severity. In critical illness: low T3 and T4 are associated with worse prognosis. T4 and TSH may rise during recovery. No treatment indicated; thyroid hormone replacement does not improve outcomes. Correct the underlying illness.

\medterm{Familial Hypocalciuric Hypercalcemia}-- Autosomal dominant condition caused by inactivating mutations in the calcium-sensing receptor gene leading to a higher set point for calcium-mediated PTH suppression. Presents with mild hypercalcemia, inappropriately normal or mildly elevated PTH, and low urinary calcium excretion with calcium-to-creatinine clearance ratio less than 0.01. Benign condition that does not require parathyroidectomy. Distinguished from primary hyperparathyroidism by urinary calcium excretion to prevent unnecessary surgery. Genetic testing confirms the diagnosis and identifies the specific mutation in the calcium-sensing receptor gene for definitive classification.

\medterm{Follicular Thyroid Carcinoma}-- Second most common thyroid carcinoma accounting for approximately ten percent of cases. Distinguished from follicular adenoma by capsular or vascular invasion on histological examination which cannot be determined by fine-needle aspiration. Metastasizes hematogenously to bone, lung, and liver. Treatment includes total thyroidectomy with radioactive iodine ablation. Prognosis is generally favorable though worse than papillary carcinoma particularly with distant metastases.

\medterm{Galactina}-- A hormone produced by the placenta that stimulates prolactin secretion from the anterior pituitary. Involved in mammary gland development during pregnancy and the onset of lactation postpartum. Elevated levels are associated with pathological galactorrhoea and certain pituitary tumours. Measurement is a component of the assessment of hyperprolactinaemia.

\medterm{Galactorrhoea}-- The spontaneous flow of milk from the breast not associated with childbirth or breastfeeding. Galactorrhoea can be unilateral or bilateral and is caused by hyperprolactinaemia (elevated prolactin). Aetiologies: pituitary prolactinoma (most common), hypothalamic-pituitary stalk compression, medications (antipsychotics, metoclopramide, verapamil, oral contraceptives, opioids), hypothyroidism (TRH stimulates prolactin release), chronic renal failure, chest wall irritation (herpes zoster, surgery, trauma). Galactorrhoea may also be idiopathic. Diagnostic workup: serum prolactin, TSH, pregnancy test (beta-hCG), MRI pituitary with contrast. Management: treat underlying cause; dopamine agonists (cabergoline, bromocriptine) for prolactinoma.

\medterm{Galactosaemia}-- A rare autosomal recessive disorder of galactose metabolism caused by deficiency of one of three enzymes: galactose-1-phosphate uridyltransferase (GALT, classic galactosaemia, type I), galactokinase (GALK, type II), or UDP-galactose epimerase (GALE, type III). Classic galactosaemia (GALT deficiency, incidence 1:30,000-60,000) presents in the neonatal period within days of milk feeding: vomiting, diarrhoea, failure to thrive, jaundice, hepatomegaly, cataracts (due to galactitol accumulation), hypoglycaemia, and E. coli sepsis. Diagnosis: newborn screening (GALT enzyme activity in dried blood spot), urine reducing substances. Long-term complications despite dietary restriction: speech and learning difficulties, ataxia, ovarian failure (in females), decreased bone density. Treatment: lifelong lactose/galactose-free diet.

\medterm{Gastrinoma}-- Gastrin-secreting neuroendocrine tumor causing Zollinger-Ellison syndrome characterized by gastric acid hypersecretion leading to severe peptic ulcer disease, diarrhea, and esophagitis. Approximately sixty percent are malignant with liver metastases. Multiple endocrine neoplasia type 1 association in approximately twenty percent. Diagnosis is by elevated fasting gastrin level and positive secretin stimulation test. Treatment includes high-dose proton pump inhibitors for acid control and surgical resection of localised disease. Somatostatin analogues control symptoms in functioning tumours. Chemotherapy and peptide receptor radionuclide therapy are options for metastatic disease.

\medterm{Gestational Diabetes Mellitus}-- Glucose intolerance first recognized during pregnancy affecting approximately six to nine percent of pregnancies. Results from insulin resistance induced by placental hormones including human placental lactogen, cortisol, and progesterone combined with insufficient beta cell compensation. Risk factors include advanced maternal age, obesity, family history of diabetes, previous gestational diabetes, polycystic ovary syndrome, and certain ethnicities. Universal screening is performed at twenty-four to twenty-eight weeks of gestation with a seventy-five gram oral glucose tolerance test. Treatment includes lifestyle modification and insulin or metformin if glycaemic targets are not met. Glucose levels normalise after delivery but affected women have increased long-term risk of type two diabetes.

\medterm{Gigantism}-- A rare condition of excessive growth and tall stature caused by growth hormone (GH) hypersecretion that begins before epiphyseal plate closure in childhood, before the onset of puberty. The most common cause is a GH-secreting pituitary adenoma (usually a benign somatotroph adenoma, often macroadenoma greater than 1 cm), which may be sporadic or associated with genetic syndromes (McCune-Albright syndrome, multiple endocrine neoplasia type 1, Carney complex, familial isolated pituitary adenoma from germline mutations in the AIP gene). Presents with accelerated linear growth, tall stature, large hands and feet, coarse facial features, and delayed skeletal maturation. Diagnosis is confirmed by elevated IGF-1 and GH levels that fail to suppress on glucose tolerance testing. Pituitary MRI identifies the responsible adenoma. Treatment includes transsphenoidal adenomectomy, medical therapy with somatostatin analogues, and rarely radiotherapy. Early diagnosis and intervention improve final height outcomes.

\medterm{Glucagonoma}-- A rare pancreatic neuroendocrine tumor secreting glucagon. Classic presentation: 4 Ds — diabetes mellitus (mild, new onset), dermatitis (necrolytic migratory erythema, characteristic painful, erythematous, blistering rash on groin, perineum, extremities), diarrhea, and deep vein thrombosis. Presents with weight loss, glossitis, and hypoaminoacidemia. Most are malignant and metastatic at diagnosis. Diagnosis: elevated glucagon level (>500 pg/mL) and pancreatic tumor on imaging. Treatment: surgical resection; somatostatin analogs for symptom control; chemotherapy and PRRT for advanced disease.

\medterm{Glucocorticoid Replacement Therapy}-- Replacement of cortisol deficiency in adrenal insufficiency using hydrocortisone in divided doses. Typical regimen is ten to fifteen milligrams in the morning and five to ten milligrams in the early afternoon mimicking the physiological cortisol diurnal rhythm. Alternative preparations include prednisolone and dexamethasone though these have longer half-lives and are less physiological. Fludrocortisone is added for mineralocorticoid replacement in primary adrenal insufficiency. Stress dosing requires doubling or tripling of the daily hydrocortisone dose during intercurrent illness, fever, or minor surgery, with parenteral administration for major surgery or when oral intake is not possible. Patient education about sick-day rules and provision of an emergency injectable glucocorticoid kit are essential to prevent adrenal crisis.

\medterm{Glucocorticoids}-- A class of steroid hormones produced by the zona fasciculata of the adrenal cortex, regulated by the hypothalamic-pituitary-adrenal (HPA) axis. Cortisol is the primary endogenous glucocorticoid, essential for: (1) carbohydrate metabolism (gluconeogenesis, glycogenolysis, peripheral insulin resistance); (2) protein metabolism (catabolism in muscle, connective tissue, bone); (3) fat metabolism (lipolysis, fat redistribution to face and trunk); (4) anti-inflammatory and immunosuppressive effects (inhibit phospholipase A2, COX-2, pro-inflammatory cytokines); (5) stress response; (6) maintenance of blood pressure (permissive effect on catecholamines). Synthetic glucocorticoids (prednisolone, dexamethasone, betamethasone, hydrocortisone) are used therapeutically for: autoimmune and inflammatory diseases (rheumatoid arthritis, asthma, inflammatory bowel disease, SLE), allergic reactions, organ transplant rejection, adrenal insufficiency (replacement therapy), cerebral oedema (dexamethasone), and fetal lung maturation (betamethasone). Side effects of chronic use: Cushing syndrome (central obesity, moon face, buffalo hump, striae), osteoporosis, diabetes, hypertension, immunodeficiency, cataract, growth suppression in children.

\medterm{Glucose Management Indicator}-- Derived metric from continuous glucose monitoring that estimates the hemoglobin A1c equivalent based on mean sensor glucose values. Formula converts mean glucose to an A1c-equivalent value allowing comparison between CGM data and traditional A1c. Provides a more comprehensive picture of glycemic control than A1c alone because it reflects glucose variability and the complete glucose profile rather than a weighted average.

\medterm{Glycaemia}-- The concentration of glucose in the blood, normally maintained within a narrow range (fasting: 3.9-5.5 mmol/L or 70-100 mg/dL) through the coordinated action of insulin (lowers glucose) and counter-regulatory hormones (glucagon, cortisol, catecholamines, growth hormone). Disorders of glycaemia: hyperglycaemia (elevated glucose, diagnostic of diabetes mellitus when fasting >7.0 mmol/L or HbA1c >48 mmol/mol), hypoglycaemia (low glucose, <3.0 mmol/L clinically significant, causes: insulin overdose, sulfonylurea excess, insulinoma, liver failure, alcohol excess). Glycaemic control is assessed by: self-monitored blood glucose (SMBG), continuous glucose monitoring (CGM), HbA1c (glycated haemoglobin, reflects average glucose over 2-3 months).

\medterm{Glycosuria}-- The presence of glucose in the urine. Normally, the renal tubules reabsorb filtered glucose via sodium-glucose co-transporters (SGLT2 in proximal tubule). Glycosuria occurs when blood glucose exceeds the renal threshold (approximately 10 mmol/L or 180 mg/dL), overwhelming reabsorptive capacity. Causes: (1) hyperglycaemia (diabetes mellitus, most common); (2) renal glycosuria (benign condition with low renal threshold, due to SGLT2 mutation); (3) pregnancy (decreased renal threshold due to increased GFR). Diagnostic significance: glycosuria is a screening finding for diabetes; its absence does not exclude diabetes. SGLT2 inhibitors (dapagliflozin, empagliflozin) are a class of antidiabetic drugs that intentionally induce glycosuria by blocking renal glucose reabsorption.

\medterm{Goiter}-- Enlargement of the thyroid gland, which can be diffuse or nodular. Causes: iodine deficiency (most common worldwide), Hashimoto thyroiditis, Graves disease, multinodular goiter, and thyroid cancer. May be associated with normal, hyperthyroid, or hypothyroid function. Large goiters can cause compressive symptoms: dysphagia, dysphonia, and stridor. Diagnosis: thyroid ultrasound, TSH, and thyroid scan. Treatment depends on etiology and symptoms.

\medterm{Goitre}-- An abnormal enlargement of the thyroid gland. Goitre may be diffuse (uniform enlargement, e.g., Graves disease, Hashimoto thyroiditis, endemic iodine deficiency) or nodular (single nodule or multinodular). Functional status: toxic (hormone overproduction, hyperthyroidism) or non-toxic (euthyroid or hypothyroid). Aetiologies: iodine deficiency (most common worldwide), autoimmune (Graves disease, Hashimoto), inflammatory (subacute thyroiditis, painless thyroiditis), neoplastic (benign adenoma, thyroid cancer), congenital (dyshormonogenesis), infiltrative (amyloidosis, sarcoidosis). Symptoms: neck swelling, globus sensation, dysphagia, dyspnoea (tracheal compression), hoarseness (recurrent laryngeal nerve compression). Pemberton sign: facial flushing and venous distension with arm elevation, indicates thoracic inlet obstruction. Workup: TSH, free T4, thyroid ultrasound, radioactive iodine uptake scan. Management depends on aetiology and symptoms: levothyroxine for hypothyroidism, antithyroid drugs/radioiodine/surgery for hyperthyroidism, thyroidectomy for compressive symptoms or malignancy.

\medterm{Gonadotrophins}-- Hormones that stimulate the gonads. The anterior pituitary secretes follicle-stimulating hormone (FSH) and luteinising hormone (LH), both glycoprotein hormones sharing a common alpha subunit (identical to TSH and hCG alpha) with a unique beta subunit. FSH: females--stimulates ovarian follicle growth and oestrogen production; males--stimulates Sertoli cells for spermatogenesis. LH: females--triggers ovulation and corpus luteum formation (progesterone); males--stimulates Leydig cells to produce testosterone. Human chorionic gonadotrophin (hCG) is produced by the placenta and maintains the corpus luteum in early pregnancy. Therapeutically, gonadotrophins are used in assisted reproductive technology (IVF, ovulation induction, spermatogenesis induction).

\medterm{Gonads}-- The primary reproductive organs: testes in males and ovaries in females. Functions: (1) gamete production (spermatogenesis in testes, oogenesis in ovaries); (2) sex steroid hormone secretion (testosterone from testicular Leydig cells; oestrogen and progesterone from ovarian follicular cells and corpus luteum). Gonadal function is regulated by the hypothalamic-pituitary-gonadal axis: GnRH from hypothalamus stimulates gonadotrophins (LH and FSH) from anterior pituitary, which act on the gonads. Disorders: hypogonadism (primary: gonadal failure; secondary: pituitary/hypothalamic deficiency), hypergonadism (precocious puberty), gonadal dysgenesis (Turner syndrome, Klinefelter syndrome).

\medterm{Graves Disease}-- Autoimmune disorder and most common cause of hyperthyroidism characterized by thyroid-stimulating immunoglobulins that bind and activate the thyroid-stimulating hormone receptor causing unregulated thyroid hormone synthesis and secretion. The diffuse goiter reflects TSH receptor stimulation. Ophthalmopathy occurs in approximately twenty-five percent with proptosis, periorbital edema, diplopia, and in severe cases optic nerve compression. Pretibial myxedema is a localized dermopathy. Diagnosis is confirmed by elevated free T4 and suppressed TSH, positive thyroid-stimulating immunoglobulins, and increased radioactive iodine uptake on thyroid scintigraphy.

\medterm{Graves' Disease}-- An autoimmune disorder causing hyperthyroidism, the most common cause of hyperthyroidism (60-80% of cases). Pathophysiology: autoantibodies (thyroid-stimulating immunoglobulins, TSI) bind to and activate the TSH receptor. Incidence: 1-2 per 1000/year, 5-10 times more common in women; peak age 30-50. Clinical features: (1) hypermetabolic state--weight loss despite increased appetite, heat intolerance, tremor, palpitations, sweating, tachycardia; (2) goitre--diffuse, symmetric, may have a bruit; (3) eye disease (Graves ophthalmopathy)--lid retraction, proptosis (exophthalmos), periorbital oedema, diplopia, optic neuropathy; (4) dermopathy--pretibial myxoedema (rare); (5) thyroid acropachy (rare, clubbing). Diagnosis: low TSH, elevated free T4 and/or T3, positive TSI antibodies. Management: beta-blockers (symptom control), antithyroid drugs (methimazole/carbimazole, propylthiouracil in first trimester), radioactive iodine-131, thyroidectomy. Complications: thyroid storm (life-threatening), atrial fibrillation, osteoporosis, ophthalmopathy.

\medterm{Growth Hormone Deficiency}-- Insufficient secretion of growth hormone from the anterior pituitary, causing short stature in children and metabolic effects in adults. In children: delayed growth, short stature, delayed bone age, and increased fat mass. In adults: decreased muscle mass, increased fat mass, reduced bone density, fatigue, and dyslipidemia. Diagnosis: low IGF-1, GH stimulation testing. Treatment: recombinant human GH (somatropin) replacement.

\medterm{Growth Hormone Physiology}-- Growth hormone is secreted in a pulsatile pattern by pituitary somatotroph cells regulated by hypothalamic GHRH which stimulates release and somatostatin which inhibits release. GH acts directly on tissues and indirectly through IGF-1 produced primarily by the liver. GH promotes linear growth in children and in adults regulates body composition with anabolic effects on muscle catabolic effects on adipose tissue effects on bone metabolism and glucose homeostasis. GH secretion decreases with aging, contributing to sarcopenia and increased adiposity in the elderly. Testing for GH deficiency requires provocative testing with insulin tolerance test or glucagon stimulation test, as random GH levels are unreliable due to pulsatile secretion.

\medterm{Gynaecomastia}-- The benign enlargement of breast tissue in males. Pathophysiology: an imbalance between oestrogen (stimulates breast tissue growth) and androgen (suppresses) activity, resulting in increased oestrogen-to-testosterone ratio. Gynaecomastia can be: (1) Physiologic: neonatal (transient, maternal oestrogen), pubertal (60-70% of boys, peak age 13-14, typically resolves within 2 years), senescent (ageing-related increase in aromatase activity and decrease in testosterone). (2) Pathologic: hypogonadism (primary--Klinefelter syndrome, secondary--pituitary failure), hyperthyroidism (increased aromatisation), liver cirrhosis (impaired oestrogen metabolism), testicular tumours (Leydig cell, Sertoli cell, choriocarcinoma--hCG stimulates oestrogen), renal failure, refeeding after starvation. (3) Drug-induced: anti-androgens (spironolactone, cimetidine, finasteride, bicalutamide), oestrogens, anabolic steroids, GnRH agonists, calcium channel blockers, digoxin, isoniazid, HIV drugs, tricyclic antidepressants, methadone, marijuana. Diagnosis: true gynaecomastia (palpable concentric disc of firm tissue >0.5 cm) vs pseudogynaecomastia (adipomastia, fatty enlargement without breast tissue). Workup: breast examination, testicular exam, LFT, TFT, serum hCG, LH, FSH, testosterone (total and free), prolactin, oestradiol. Imaging: mammography if asymmetric or suspicious. Management: reassurance for physiologic gynaecomastia (spontaneous regression); discontinuation of offending drugs; medical therapy: tamoxifen (SERM, reduces oestrogen effect on breast) or raloxifene; severe/long-standing cases (fibrotic stage >12 months): subcutaneous mastectomy (liposuction or surgical excision).

\medterm{Hashimoto Thyroiditis}-- Autoimmune thyroiditis characterized by lymphocytic infiltration and progressive destruction of thyroid follicles leading to hypothyroidism. Most common cause of hypothyroidism in iodine-sufficient areas. Anti-thyroid peroxidase and anti-thyroglobulin antibodies are present. The thyroid may be diffusely enlarged and firm or atrophic. Patients present with fatigue, weight gain, cold intolerance, constipation, dry skin, and depression. Subclinical hypothyroidism with elevated TSH and normal free T4 levels is common and may progress to overt hypothyroidism. Treatment is levothyroxine replacement titrated to achieve normal TSH levels.

\medterm{Healthy Eating for Diabetes}-- A dietary approach designed to maintain stable blood glucose levels in individuals with diabetes mellitus. It emphasizes carbohydrate consistency, high fiber intake, lean proteins, and healthy fats while minimizing refined sugars and simple carbohydrates. Meal planning often involves carbohydrate counting or the plate method, and coordination with medication and physical activity is essential.

\medterm{Hemoglobin A1c}-- Glycated hemoglobin reflecting average blood glucose concentration over the preceding two to three months corresponding to the red blood cell lifespan. Formed by non-enzymatic glycation of hemoglobin at the N-terminal valine of the beta chain. The percent of glycated hemoglobin is directly proportional to mean glucose levels. Used for diagnosis of diabetes with greater than or equal to 6.5 percent indicating diabetes, 5.7 to 6.4 percent indicating prediabetes, and less than 5.7 percent indicating normoglycaemia. Limitations include interference from haemoglobinopathies, anaemia, and conditions affecting red blood cell lifespan, which may give falsely low or high values independent of glycaemic control.

\medterm{Hemoglobin A1c Variability}-- Factors that affect hemoglobin A1c accuracy include hemoglobinopathies such as sickle cell disease and thalassemias which alter red blood cell lifespan and glycation rates, iron deficiency anemia which shortens red blood cell survival lowering A1c, vitamin B12 and folate deficiency, pregnancy, chronic kidney disease, erythropoietin therapy, and recent blood transfusion. Conditions that shorten red blood cell lifespan lower A1c while those that prolong lifespan raise A1c independent of glycemic control. When A1c and glucose data are discordant, alternative methods such as fructosamine or glycated albumin may provide useful complementary information, and a careful review for conditions affecting red blood cell turnover is warranted.

\medterm{Hirsutism}-- Excess terminal hair growth in women in a male-pattern distribution (face--upper lip, chin, sideburns; chest, back, lower abdomen, inner thighs). Hirsutism affects 5-10% of women of reproductive age. Pathophysiology: increased androgen production or increased end-organ sensitivity. The modified Ferriman-Gallwey (mFG) score quantifies hirsutism (score >8 indicates hirsutism). Causes: (1) Polycystic ovary syndrome (PCOS, 70-80% of cases): elevated LH:FSH ratio, hyperandrogenism, anovulation. (2) Idiopathic hirsutism (15-20%): normal serum androgens, increased 5-alpha-reductase activity (converts testosterone to DHT in the hair follicle). (3) Non-classic congenital adrenal hyperplasia (NCCAH, 5-10%): 21-hydroxylase deficiency (elevated 17-hydroxyprogesterone). (4) Androgen-secreting tumours (ovarian or adrenal, rare, rapidly progressive): testosterone >200 ng/dL, DHEAS >700 mcg/dL. (5) Cushing syndrome: central obesity, striae, hypertension, hyperglycaemia. (6) Medications: danazol, testosterone, anabolic steroids, ciclosporin, valproate, phenytoin. (7) Other: HAIR-AN syndrome (hyperandrogenism, insulin resistance, acanthosis nigricans), hyperprolactinaemia. Diagnosis: serum total and free testosterone, SHBG, DHEAS, 17-OHP, LH, FSH, prolactin, TSH, free T4; pelvic ultrasound for PCOS. Management: (1) Mechanical: hair removal by shaving, waxing, electrolysis, laser hair removal (most effective, targets melanin in hair bulb). (2) Pharmacological: combined oral contraceptive pill (first line for PCOS, reduces ovarian androgen production, increases SHBG), anti-androgens (spironolactone 100-200 mg/day, finasteride, cyproterone acetate, flutamide--teratogenic, requires contraception), metformin (insulin sensitizer, modest benefit), eflornithine cream (topical, inhibits ornithine decarboxylase, slows hair growth). (3) Treatment of underlying cause (CAH: glucocorticoids; tumour: surgical removal).

\medterm{Hormone}-- A chemical messenger produced by endocrine glands, secreted into the bloodstream, and transported to distant target cells where it binds to specific receptors to exert a physiological effect. Hormones regulate: growth and development (growth hormone, thyroid hormone, sex steroids), metabolism (insulin, glucagon, cortisol, thyroid hormone, adrenaline), reproduction (LH, FSH, oestrogen, progesterone, testosterone, prolactin), fluid and electrolyte balance (ADH, aldosterone), calcium homeostasis (PTH, calcitonin, vitamin D), and stress response (cortisol, adrenaline, noradrenaline). Hormone classification: (1) Peptide/protein hormones (water-soluble): hypothalamic releasing factors, pituitary hormones, insulin, glucagon, PTH, calcitonin; bind to cell surface receptors (GPCRs, receptor tyrosine kinases), act via second messengers (cAMP, IP3/DAG, Ca2+). (2) Steroid hormones (lipid-soluble): cortisol, aldosterone, oestrogen, progesterone, testosterone, vitamin D; synthesised from cholesterol, bind to intracellular receptors (cytoplasmic/nuclear receptors), act as transcription factors (slow genomic effects, hours-days). (3) Amino acid-derived: thyroid hormones (T3/T4, iodothyronines, nuclear receptors), catecholamines (adrenaline/noradrenaline, GPCRs), melatonin. (4) Prostaglandins/leukotrienes (eicosanoids): paracrine/autocrine. Hormone secretion is regulated by feedback loops (most commonly negative feedback: e.g., T4 inhibits TRH and TSH; insulin lowers glucose which reduces further insulin secretion) and by the hypothalamic-pituitary axis. Disorders: overproduction (hyperfunction) or underproduction (hypofunction) of hormones cause endocrine diseases.

\medterm{Hormone Replacement Therapy (HRT)}-- The administration of exogenous oestrogen (with or without progestogen) to relieve symptoms of oestrogen deficiency and prevent long-term health consequences in postmenopausal women. Indications: moderate-severe vasomotor symptoms (hot flushes, night sweats, 80% of women experience at some point), urogenital atrophy (vaginal dryness, dyspareunia, recurrent UTIs), prevention of osteoporosis (in women at high fracture risk who cannot tolerate bisphosphonates), and premature ovarian insufficiency (POI, <40 years, HRT needed at least until average age of menopause). Estrogen-only: for women post-hysterectomy; reduces risk of endometrial cancer. Combined oestrogen-progestogen: for women with intact uterus (unopposed oestrogen increases endometrial cancer risk; progestogen protects by inducing endometrial shedding). Administration routes: oral (conjugated equine oestrogens 0.3-0.625 mg, oestradiol 1-2 mg; medroxyprogesterone acetate 2.5-5 mg), transdermal (patches/gel, avoids first-pass hepatic metabolism, lower VTE risk), vaginal (cream, ring, tablets for urogenital symptoms only, minimal systemic absorption). Risks: VTE (2-4 fold increased, highest in first year), breast cancer (5-6 extra cases per 10,000 women-years with combination HRT >5 years, lower with oestrogen-only), stroke, gallbladder disease, endometrial cancer (if unopposed oestrogen). The Women's Health Initiative (WHI) 2002 found overall risks outweighed benefits for chronic disease prevention in older women (mean age 63), leading to reduced HRT use. Current guidelines (NICE, IMS): HRT is safe and effective for symptom management in women aged <60 or <10 years from menopause onset; lowest effective dose for shortest duration (no arbitrary upper limit for symptom management); tibolone (synthetic steroid) is an alternative. Contraindications: breast cancer (current or prior), oestrogen-dependent tumours, unexplained vaginal bleeding, active VTE, history of VTE if not anticoagulated, arterial disease, pregnancy, porphyria.

\medterm{Hungry Bone Syndrome}-- Severe hypocalcemia occurring after parathyroidectomy particularly in patients with severe preoperative hyperparathyroidism and high bone turnover. Rapid bone remineralization after removal of excess PTH sequesters calcium and phosphorus from the serum causing symptomatic hypocalcemia. Risk factors include severe preoperative hyperparathyroidism, high alkaline phosphatosis, and osteitis fibrosa cystica. Prevention includes preoperative calcium and vitamin D supplementation. Treatment is aggressive intravenous and oral calcium and calcitriol supplementation, often requiring high doses that are gradually tapered as bone turnover normalises over weeks to months.

\medterm{Hyperaldosteronism Complications}-- Beyond hypertension and hypokalemia, primary hyperaldosteronism causes target organ damage including left ventricular hypertrophy, fibrosis, stroke, coronary artery disease, atrial fibrillation, and chronic kidney disease. These complications occur at higher rates than in essential hypertension of similar severity. Cardiovascular risk persists even after treatment of the hyperaldosteronism supporting early diagnosis and intervention. Screening is recommended for all patients with resistant hypertension, spontaneous or diuretic-induced hypokalaemia, or adrenal incidentaloma, and is cost-effective given the high prevalence and cardiovascular morbidity.

\medterm{Hypercalcemia}-- Serum calcium greater than 10.5 milligrams per deciliter. Causes summarized by mnemonic mnemonic MALIGNANT: Myeloma, Metastases, Metabolic bone disease, Hyperparathyroidism, Adrenal insufficiency, Paget disease, Immobilization, Thiazide diuretics, Addison disease, Neonatal, Tumor-derived PTHrP. Symptoms include stones bones groans and psychic moans referring to nephrolithiasis, bone pain, constipation, and depression. Treatment includes saline hydration, loop diuretics after hydration, bisphosphonates for sustained calcium lowering, and denosumab for refractory hypercalcaemia. Underlying cause should be addressed.

\medterm{Hyperglycemia}-- Elevated blood glucose levels above normal range (>140 mg/dL postprandial or >126 mg/dL fasting). Acute hyperglycemia can cause osmotic diuresis, dehydration, and electrolyte disturbances. Chronic hyperglycemia leads to microvascular (retinopathy, nephropathy, neuropathy) and macrovascular complications. Management: treat underlying diabetes, insulin or oral hypoglycemics as indicated, and address precipitating factors.

\medterm{Hypernatremia}-- Serum sodium greater than 145 mEq/L. Always indicates free water deficit relative to sodium. Causes include inadequate water intake particularly in elderly and cognitively impaired patients, diabetes insipidus, osmotic diuresis from hyperglycemia, and GI or renal water losses. Treatment is free water replacement with rate of correction dependent on chronicity. Overly rapid correction risks cerebral edema particularly in chronic hypernatremia. D5W or hypotonic saline is used for correction.

\medterm{Hyperosmolar Hyperglycemic State}-- Hyperglycemic emergency characterized by severe hyperglycemia greater than 600 milligrams per deciliter, hyperosmolality greater than 320 milliosmoles per kilogram, and profound dehydration without significant ketosis. Primarily occurs in type two diabetes where enough endogenous insulin is present to prevent ketogenesis but insufficient to control glucose. Presents with extreme dehydration, altered mental status, seizures, and focal neurological deficits. Mortality rates range from five to twenty percent. Treatment is aggressive IV fluid resuscitation, insulin infusion to lower glucose and shift fluids intracellularly, and careful electrolyte monitoring and replacement.

\medterm{Hyperparathyroidism}-- A disorder of excessive parathyroid hormone (PTH) secretion from one or more of the four parathyroid glands, classified as primary, secondary, or tertiary. Primary hyperparathyroidism is caused by a single parathyroid adenoma (80-85 percent), multigland hyperplasia (10-15 percent, associated with MEN1 and MEN2A), or parathyroid carcinoma (less than 1 percent). It is the most common cause of hypercalcemia in the outpatient setting, with a prevalence of 1-4 per 1000 adults, most common in postmenopausal women with a prevalence of approximately two percent in those over fifty-five years. Symptoms of hypercalcaemia include polyuria, polydipsia, nephrolithiasis, bone pain, constipation, fatigue, and cognitive dysfunction. Severe hypercalcaemia above fourteen milligrams per decilitre may cause altered mental status, cardiac arrhythmias, and acute kidney injury. Diagnosis requires elevated serum calcium with inappropriately normal or elevated PTH. Localisation is performed with sestamibi scan, ultrasound, or four-dimensional CT to identify the abnormal parathyroid gland. Parathyroidectomy is curative for symptomatic or complicated disease. Asymptomatic patients may be monitored if they meet specific criteria including age, calcium level, bone density, and renal function parameters.

\medterm{Hyperprolactinemia}-- Elevated serum prolactin levels most commonly caused by prolactin-secreting pituitary adenomas termed prolactinomas. Other causes include medications particularly antipsychotics and metoclopramide, hypothyroidism, stalk effect from non-functioning pituitary adenomas, renal failure, and chest wall lesions. Women present with galactorrhea, amenorrhea, and infertility. Men present with decreased libido, erectile dysfunction, and rarely galactorrhea. Macroprolactinomas may present with mass effects including headache, visual field defects, and hypopituitarism. Diagnosis includes serum prolactin measurement to exclude stalk effect, MRI of the pituitary, and full pituitary hormone assessment to evaluate for hypopituitarism. Treatment of prolactinomas is with dopamine agonists, while non-functioning adenomas causing mass effects are managed with transsphenoidal surgery.

\medterm{Hyperthyroidism}-- A condition of excess thyroid hormone production and secretion from the thyroid gland, resulting in thyrotoxicosis (elevated serum T4 and T3 concentrations). The most common cause is Graves disease (60-80 percent of cases), an autoimmune disorder caused by thyroid-stimulating immunoglobulins (TSI) that bind to and activate the TSH receptor on thyroid follicular cells, leading to unregulated thyroid hormone production and thyroid gland enlargement. Other causes include toxic multinodular goiter (toxic nodular goiter) caused by autonomously functioning thyroid nodules in older adults with long-standing multinodular goiter, often presenting with subclinical hyperthyroidism, atrial fibrillation, and cardiovascular complications. Other causes include thyroiditis (subacute, silent, postpartum), exogenous thyroid hormone excess, and rare causes such as struma ovarii, TSH-secreting pituitary adenoma, and hCG-mediated hyperthyroidism in trophoblastic disease. Clinical features reflect increased metabolic rate: weight loss despite increased appetite, heat intolerance, diaphoresis, tremor, palpitations, tachycardia, atrial fibrillation, anxiety, irritability, insomnia, frequent bowel movements, oligomenorrhoea, and proximal myopathy. In older patients, apathetic hyperthyroidism may present with lethargy, depression, weight loss, and atrial fibrillation without classic signs of adrenergic excess. Laboratory findings include suppressed TSH with elevated free T4 and T3. T3 toxicosis, where only T3 is elevated with normal T4, may occur early in Graves disease or in toxic nodular goiter. Thyroid scintigraphy shows diffuse increased uptake in Graves disease, heterogeneous uptake in toxic multinodular goiter, and a focal hot nodule in toxic adenoma, whereas uptake is low in thyroiditis.

\medterm{Hyperthyroidism (Thyrotoxicosis)}-- See Hyperthyroidism.

\medterm{Hypocalcemia}-- Serum calcium less than 8.5 milligrams per deciliter or ionized calcium less than 4.5 milligrams per deciliter. Causes include hypoparathyroidism, vitamin D deficiency, hypomagnesemia, hyperphosphatemia, and critical illness. Symptoms include perioral numbness, paresthesias, muscle cramps, carpopedal spasm, Chvostek sign, Trousseau sign, seizures, and QTc prolongation. Acute symptomatic hypocalcemia is treated with IV calcium gluconate. Chronic management includes oral calcium and calcitriol supplementation, along with magnesium repletion if hypomagnesaemia is present to correct the impaired PTH secretion.

\medterm{Hypoglycaemia}-- An abnormally low blood glucose level, defined as plasma glucose <3.0 mmol/L (<54 mg/dL, the threshold for neuroglycopenic symptoms). Whipple triad confirms hypoglycaemia: (1) symptoms consistent with hypoglycaemia, (2) low plasma glucose at the time of symptoms, (3) relief of symptoms after glucose administration. Symptoms: (1) Autonomic (adrenergic): tremor, palpitations, sweating, anxiety, hunger, paraesthesiae (mediated by adrenaline and noradrenaline, occur at glucose <3.8 mmol/L). (2) Neuroglycopenic: confusion, difficulty concentrating, drowsiness, behavioural changes, slurred speech, ataxia, visual disturbances, seizures, loss of consciousness, coma (glucose <2.8 mmol/L). Causes in diabetes: excessive insulin or sulfonylurea dose, missed meals, unplanned exercise, alcohol consumption, illness. Causes in non-diabetics: insulinoma (pancreatic beta cell tumour, fasting hypoglycaemia), reactive (postprandial) hypoglycaemia (alimentary--post-gastrectomy dumping; idiopathic), insulin autoimmune syndrome (Hirata disease, insulin autoantibodies), factitious (insulin or sulfonylurea abuse), critical illness (liver failure, renal failure, sepsis, heart failure), medications (salicylates, quinolones, pentamidine), alcohol (inhibits gluconeogenesis), hormone deficiencies (cortisol, growth hormone, glucagon), and inborn errors of metabolism (glycogen storage diseases, fatty acid oxidation defects). Management: (1) Conscious, able to swallow: oral fast-acting carbohydrate (15-20 g glucose, e.g., 150-200 mL fruit juice, 3-4 glucose tablets, 5-6 sugar lumps, or 4 teaspoons of sugar). Repeat blood glucose testing after 15 minutes; give additional 15-20 g if still low. Once stabilised, give a snack with complex carbohydrate to prevent recurrence. (2) Unconscious/unable to swallow: IM glucagon (1 mg for adults, 0.5 mg for children; or SC, but IM is more reliable); IV dextrose 50% (50 mL, 25 g glucose) over 1-2 minutes via large vein (extravasation causes severe tissue necrosis). Check glucose 10-15 minutes later; repeat if needed. Start IV dextrose 10% infusion if no IV access maintained. Prevent recurrence by adjusting diabetes medications and addressing underlying cause.

\medterm{Hypoglycemia}-- Low blood glucose (<70 mg/dL) with autonomic (sweating, palpitations, tremor, anxiety) and neuroglycopenic (confusion, seizures, coma) symptoms. In diabetic patients, common causes: excessive insulin or sulfonylurea doses, missed meals, exercise, and alcohol intake. Severe hypoglycemia requires glucagon or IV dextrose. Whipple triad: symptoms consistent with hypoglycemia, low plasma glucose, and symptom resolution with glucose administration.

\medterm{Hypoglycemia in Diabetes}-- Blood glucose less than 70 milligrams per deciliter in patients with diabetes. Classification by level: level one less than 70, level two less than 54 which is clinically significant hypoglycemia, and level three severe hypoglycemia requiring assistance. Causes include insulin or sulfonylurea overdose, missed meals, excessive exercise, and alcohol consumption. Counter-regulatory hormone responses including glucagon and epinephrine cause tremor, sweating, palpitations, and anxiety. Neuroglycopenic symptoms including confusion, drowsiness, visual disturbance, slurred speech, seizures, and loss of consciousness occur with progressive cerebral glucose deprivation. Treatment is immediate ingestion of fast-acting carbohydrate for conscious patients, or intramuscular glucagon or intravenous dextrose for severe hypoglycaemia with altered consciousness.

\medterm{Hypoglycemia Unawareness}-- Loss of the warning symptoms of hypoglycemia including tremor, sweating, and palpitations occurring in patients with recurrent hypoglycemia particularly type one diabetes. Results from attenuated counter-regulatory hormone responses particularly impaired glucagon and epinephrine responses. Significantly increases the risk of severe hypoglycemia. Associated with long diabetes duration, recurrent hypoglycemic episodes, and autonomic neuropathy. Management includes liberalization of glycemic targets, structured education in hypoglycaemia recognition, avoidance of hypoglycaemia to restore symptomatic awareness, and use of continuous glucose monitoring with alerts.

\medterm{Hypogonadism}-- Decreased functional activity of the gonads (testes or ovaries) with reduced sex hormone production. Male hypogonadism: low testosterone causing decreased libido, erectile dysfunction, infertility, reduced muscle mass, and osteoporosis. Female: estrogen deficiency causing amenorrhea, hot flashes, and osteoporosis. Primary: gonadal failure (e.g., Klinefelter, Turner). Secondary: pituitary/hypothalamic dysfunction (e.g., prolactinoma, hypothalamic amenorrhea). Diagnosis: sex hormone levels, LH/FSH. Treatment: hormone replacement therapy.

\medterm{Hypogonadism in Men}-- Testosterone deficiency resulting from inadequate testicular function classified as primary testicular failure or secondary hypothalamic-pituitary dysfunction. Primary causes include Klinefelter syndrome, cryptorchidism, orchitis, trauma, and chemotherapy. Secondary causes include Kallmann syndrome, pituitary tumors, hyperprolactinemia, and chronic illness. Symptoms include decreased libido, erectile dysfunction, fatigue, decreased muscle mass, increased body fat, osteoporosis, and depression. Diagnosis includes morning total testosterone on at least two occasions, LH and FSH to distinguish primary from secondary hypogonadism, prolactin, and iron studies to exclude haemochromatosis. Testosterone replacement therapy is indicated for symptomatic men with confirmed deficiency after excluding contraindications including prostate cancer and untreated sleep apnoea.

\medterm{Hypogonadism in Women}-- Testosterone deficiency in women causing decreased libido, fatigue, depressed mood, and reduced sense of well-being. Causes include oophorectomy, premature ovarian insufficiency, adrenal insufficiency, hypopituitarism, and medications. Diagnosis is challenging due to the lack of well-established reference ranges for female testosterone levels. Treatment with transdermal testosterone at physiological doses may improve symptoms. Benefits and risks of long-term therapy are still under investigation, and long-term safety data are limited. Individualised treatment decisions should weigh potential benefits against possible adverse effects on cardiovascular health and breast cancer risk.

\medterm{Hypomagnesemia}-- Serum magnesium less than 1.7 milligrams per deciliter. Causes include gastrointestinal losses from diarrhea, renal losses from diuretics and alcohol, malnutrition, and medications including proton pump inhibitors. Manifestations include muscle cramps, tremor, seizures, cardiac arrhythmias including torsades de pointes, and refractory hypokalemia and hypocalcemia. Treatment is magnesium replacement orally for mild cases and intravenously for severe or symptomatic hypomagnesemia. Always consider checking serum magnesium in any patient with unexplained hypokalaemia or hypocalcaemia, as hypomagnesaemia impairs renal potassium conservation and PTH secretion, making electrolyte repletion difficult until magnesium is normalised.

\medterm{Hyponatremia}-- Serum sodium less than 135 mEq/L. Most common electrolyte disorder in hospitalized patients. Causes are classified by volume status: hypovolemic from GI losses and diuretics, euvolemic from SIADH and hypothyroidism, and hypervolemic from heart failure, cirrhosis, and nephrotic syndrome. Acute hyponatremia with neurological symptoms requires rapid correction with hypertonic saline. Chronic hyponatremia should be corrected slowly at no more than eight milliequivalents per liter per twenty-four hours to prevent osmotic demyelination syndrome. Use of vasopressin receptor antagonists (vaptans) may be considered in euvolemic and hypervolemic hyponatremia when conventional therapy is insufficient.

\medterm{Hypoparathyroidism}-- Deficiency or absence of parathyroid hormone causing hypocalcemia and hyperphosphatemia. Most commonly occurs after thyroid or parathyroid surgery from inadvertent removal of parathyroid glands. Autoimmune hypoparathyroidism may occur alone or as part of autoimmune polyendocrine syndrome type one. Symptoms of hypocalcemia include perioral tingling, muscle cramps, carpopedal spasm, Chvostek sign, Trousseau sign, seizures, and cardiac arrhythmias. Chronic hypocalcemia causes basal ganglia calcifications. Treatment requires oral calcium and active vitamin D (calcitriol) supplementation to maintain serum calcium in the low-normal range and avoid hypercalciuria.

\medterm{Hypoparathyroidism Management}-- Treatment of chronic hypocalcemia from hypoparathyroidism includes oral calcium supplementation in divided doses and active vitamin D calcitriol to enhance calcium absorption. Goal is to maintain serum calcium in the low-normal range to prevent hypercalciuria and renal complications. Thiazide diuretics reduce urinary calcium excretion. Magnesium supplementation may be needed if hypomagnesemia is present. PTH replacement therapy with recombinant human PTH is approved for refractory cases not controlled with conventional therapy, but its use is limited by cost and the need for twice-daily injections.

\medterm{Hypopituitarism}-- Deficiency of one or more pituitary hormones most commonly from pituitary adenomas, craniopharyngiomas, pituitary surgery or radiation, traumatic brain injury, or infiltrative diseases. Growth hormone deficiency is typically the first hormone lost followed by gonadotropins then TSH and ACTH. Diagnosis requires dynamic testing of pituitary reserve. Treatment is replacement of deficient hormones with levothyroxine, hydrocortisone, sex steroids, desmopressin for DI, and growth hormone.

\medterm{Hypothalamic-Pituitary-Adrenal Axis Physiology}-- The hypothalamus releases CRH in a diurnal rhythm peaking in the early morning stimulating pituitary corticotrophs to secrete ACTH which stimulates the adrenal cortex to produce cortisol. Cortisol exerts negative feedback on both the hypothalamus and pituitary suppressing CRH and ACTH release. The system is activated during physiological stress overriding normal feedback mechanisms. Assessment of axis integrity uses the overnight dexamethasone suppression test for Cushing screening and the ACTH stimulation test for adrenal insufficiency assessment.

\medterm{Hypothalamic-Pituitary-Gonadal Axis Physiology}-- The hypothalamus releases GnRH in a pulsatile fashion stimulating pituitary gonadotrophs to secrete LH and FSH. In males LH stimulates Leydig cells to produce testosterone while FSH stimulates Sertoli cells to support spermatogenesis. In females FSH stimulates follicular development and estrogen production while the LH surge triggers ovulation. Sex steroids exert negative feedback on the hypothalamus and pituitary. Disruption at any level causes reproductive dysfunction.

\medterm{Hypothalamic-Pituitary-Thyroid Axis Physiology}-- The hypothalamus secretes TRH stimulating pituitary thyrotrophs to release TSH which stimulates the thyroid gland to produce T4 and T3. T3 the more biologically active form is produced by peripheral conversion of T4 by deiodinase enzymes. T3 and T4 exert negative feedback on both the hypothalamus and pituitary suppressing TRH and TSH release. The system maintains thyroid hormone levels within a narrow range with TSH being the most sensitive indicator of axis dysfunction.

\medterm{Hypothyroidism}-- A common condition of inadequate thyroid hormone production, with a prevalence of 4-10 percent in the general population and up to 20 percent in women over 60 years. Primary hypothyroidism accounts for 99 percent of cases, most commonly caused by chronic autoimmune (Hashimoto) thyroiditis (lymphocytic infiltration leading to gradual follicular destruction, associated with elevated anti-thyroid peroxidase [TPO] and anti-thyroglobulin antibodies), followed by iatrogenic causes (radioactive iodine therapy for Graves disease, radiation therapy for head and neck cancers, and medications (amiodarone, lithium, tyrosine kinase inhibitors). Clinical features reflect reduced metabolic rate including fatigue, lethargy, cold intolerance, weight gain, constipation, dry skin, hair loss, hoarseness, bradycardia, and delayed tendon reflexes. Severe hypothyroidism may cause myxedema, pericardial effusion, and cognitive impairment. Myxedema coma is a life-threatening complication with hypothermia, altered mental status, and respiratory depression. Diagnosis is confirmed by elevated TSH with low free T4. Positive TPO antibodies indicate autoimmune aetiology. Treatment is levothyroxine replacement with a starting dose based on patient age, weight, and cardiovascular risk, titrated to achieve normal TSH. In elderly patients or those with ischaemic heart disease, a low starting dose of twenty-five to fifty micrograms daily with gradual titration is recommended to avoid precipitating myocardial ischaemia or arrhythmia.

\medterm{Inhaled Insulin}-- Rapid-acting insulin delivered via inhalation as a powder using the Afrezza device. Approved for mealtime insulin coverage in adults with diabetes. Onset of action within twelve minutes and duration of approximately three hours providing a more physiological postprandial insulin profile than subcutaneous rapid-acting analogs. Contraindicated in smokers and patients with chronic lung disease. Spirometry monitoring is required before initiation and periodically during therapy. Reduces postprandial glucose excursions and may reduce hypoglycaemia compared with injectable mealtime insulin in selected patients. Absorption variability and limited dosing options restrict widespread use.

\medterm{Insulin Injection Technique}-- Proper technique for subcutaneous insulin injection includes selecting appropriate injection sites rotating between abdomen thighs buttocks and upper arms using appropriate needle length four to eight millimeters for most adults pinching skin for thin patients to avoid intramuscular injection rotating injection sites within an area to prevent lipohypertrophy and using a new pen needle or syringe for each injection. Site selection affects absorption rate with abdomen being fastest followed by arm and thigh being slower. Proper disposal of sharps using a sharps container is essential for safety.

\medterm{Insulin Pen Devices}-- Pre-filled or reusable insulin pen devices providing convenient and accurate insulin delivery. Benefits include portability precision dosing discrete injection and reduced needle phobia. Modern pen devices provide half-unit dosing capabilities memory features and Bluetooth connectivity for integration with continuous glucose monitoring systems and insulin dose tracking applications. Insulin pen needles are typically four to eight millimeters in length with thirty-one gauge diameter for minimal discomfort. Smart insulin pens integrate with digital health platforms for automated dose logging and decision support, enhancing diabetes self-management.

\medterm{Insulin Preparations}-- Hormone synthesized by pancreatic beta cells that regulates glucose homeostasis by promoting glucose uptake in skeletal muscle and adipose tissue, suppressing hepatic glucose production, and inhibiting lipolysis. Available preparations include rapid-acting insulin analogs including lispro, aspart, and glulisine with onset fifteen to thirty minutes and duration three to five hours, short-acting regular insulin with onset thirty to sixty minutes and duration six to eight hours, intermediate-acting NPH insulin with onset one to two hours and peak four to twelve hours, long-acting insulin analogs including glargine, detemir, and degludec with relatively flat pharmacokinetic profiles and duration up to twenty-four to forty-two hours respectively, and pre-mixed preparations containing fixed ratios of rapid-acting and intermediate-acting insulin.

\medterm{Insulin Pump Therapy}-- Continuous subcutaneous insulin infusion delivering rapid-acting insulin analog through a catheter placed subcutaneously typically in the abdomen. Provides basal insulin at programmed rates and bolus doses at mealtimes. Advantages include reduced glycemic variability, improved time in range, flexibility in meal timing, and reduced hypoglycemia compared to multiple daily injections. Advanced systems integrate with continuous glucose monitors as hybrid closed-loop or automated insulin delivery systems. Disadvantages include risk of infusion site infection, catheter occlusion, and diabetic ketoacidosis from infusion set failure. Candidates require appropriate training and commitment to pump management.

\medterm{Insulin Resistance}-- Diminished biological response to physiological insulin concentrations resulting in impaired glucose disposal in peripheral tissues and failure to suppress hepatic glucose production. Pathophysiological hallmark of type two diabetes and metabolic syndrome. Develops from complex interactions between genetic susceptibility, obesity particularly visceral adiposity, physical inactivity, and inflammatory mediators. Mechanisms include post-receptor signaling defects in the insulin signaling cascade including IRS-1 phosphorylation, PI3K activation, and GLUT4 translocation to the cell membrane. Management primarily involves lifestyle modification with weight loss, physical activity, and dietary changes, along with pharmacotherapy including metformin, thiazolidinediones, and other insulin sensitizers.

\medterm{Insulin Resistance Syndromes}-- Rare inherited severe insulin resistance syndromes include Donohue syndrome also known as leprechaunism, Rabson-Mendenhall syndrome, and type A insulin resistance syndrome caused by mutations in the insulin receptor gene. Donohue syndrome is the most severe with death in infancy. Rabson-Mendenhall syndrome presents with dental and nail abnormalities and pineal hyperplasia. Type A insulin resistance presents with acanthosis nigricans, hyperandrogenism in females, and severe insulin resistance. Treatment includes high-dose insulin therapy, metformin, and consideration of recombinant human IGF-1 in selected cases. Genetic counselling is recommended for affected families.

\medterm{Kallmann Syndrome}-- Congenital form of hypogonadotropic hypogonadism caused by failure of GnRH neuron migration from the olfactory placode to the hypothalamus during embryonic development. Characterized by delayed puberty with low LH and FSH and anosmia or hyposmia from olfactory bulb hypoplasia. Associated features include renal agenesis, cleft palate, and synkinesia. Treatment is pulsatile GnRH therapy for fertility or sex steroid replacement for virilization or feminization.

\medterm{Klinefelter Syndrome}-- Most common sex chromosome aneuploidy in males with a karyotype of 47,XXY occurring in approximately one in six hundred live male births. Characterized by small firm testes, gynecomastia, tall stature with long extremities, and infertility from azoospermia. Testosterone levels are low with elevated gonadotropins reflecting primary testicular failure. Learning disabilities and behavioral issues may occur. Testosterone replacement therapy beginning at puberty improves virilization, bone density, and quality of life. Fertility options including assisted reproductive techniques should be discussed as spermatogenesis is not restored by testosterone therapy alone.

\medterm{Liddle Syndrome Treatment}-- Liddle syndrome is an autosomal dominant cause of early-onset hypertension with hypokalemia and metabolic alkalosis caused by gain-of-function mutations in ENaC sodium channel subunits. Aldosterone levels are suppressed. Unlike primary hyperaldosteronism the hypertension responds to ENaC blockers amiloride or triamterene but not to spironolactone because the channel is constitutively active independent of aldosterone. Genetic testing confirms the diagnosis. Dietary sodium restriction is important for management as it distinguishes the condition from primary hyperaldosteronism and directs appropriate therapy.

\medterm{Lipodystrophy Syndromes}-- Rare conditions characterized by selective loss of adipose tissue leading to severe insulin resistance, diabetes, and hypertriglyceridemia. Congenital generalized lipodystrophy involves loss of all body fat from birth. Familial partial lipodystrophy involves loss of limb and truncal fat with fat accumulation in face and neck. Acquired forms occur with antiretroviral therapy particularly protease inhibitors. Diagnosis is by clinical examination and imaging showing absent or abnormal fat distribution on dual-energy X-ray absorptiometry. Management focuses on metabolic complications including metformin for diabetes, fibrates and omega-3 fatty acids for hypertriglyceridaemia, and leptin replacement therapy with metreleptin for congenital generalized lipodystrophy.

\medterm{Low Testosterone (Hypogonadism)}-- Low testosterone (male hypogonadism) is a clinical syndrome of deficient testosterone production by the testes (primary) or inadequate stimulation by the pituitary/hypothalamus (secondary), defined as total testosterone <300 ng/dL on two morning measurements with consistent symptoms. Symptoms include low libido, erectile dysfunction, fatigue, depressed mood, decreased muscle mass and bone density, increased body fat, hot flashes, and reduced energy. Causes: primary (Klinefelter syndrome, testicular trauma/orchitis, chemotherapy, aging), secondary (pituitary tumors, hemochromatosis, obesity, sleep apnea, opioids, glucocorticoids). Diagnosis requires confirmatory labs: total/free testosterone, LH, FSH, prolactin, iron studies; pituitary MRI if secondary. Testosterone replacement therapy (gels, injections, pellets, buccal tablets) improves symptoms but requires monitoring of hematocrit, PSA, and fertility impact; contraindicated in untreated prostate or breast cancer.

\medterm{Maturity Onset Diabetes of the Young}-- A monogenic form of diabetes with autosomal dominant inheritance, typically presenting before age 25 with non-insulin-dependent hyperglycemia. Caused by mutations in transcription factors or enzymes of beta-cell function. MODY subtypes: GCK-MODY (MODY 2, mild fasting hyperglycemia, no treatment needed), HNF1A-MODY (MODY 3, most common, sulfonylurea-responsive), HNF4A-MODY (MODY 1, sulfonylurea-responsive). Distinguished from type 1 diabetes by negative autoantibodies and detectable C-peptide. Distinguished from type 2 diabetes by young age, lack of obesity, and strong family history. Diagnosis: genetic testing. Treatment: varies by subtype; diet, sulfonylureas, or insulin.

\medterm{McCune-Albright Syndrome}-- Sporadic condition caused by activating mutations in the GNAS1 gene encoding the Gs alpha subunit leading to constitutive activation of adenylyl cyclase. Classic triad includes polyostotic fibrous dysplasia of bone, cafe-au-lait skin spots with irregular coast-of-Maine borders, and autonomous endocrine hyperfunction including precocious puberty, hyperthyroidism, growth hormone excess, and Cushing syndrome. Diagnosis is clinical and confirmed by genetic testing of affected tissue. Treatment targets the specific endocrine abnormality with surgery, medical therapy, or radiation as appropriate for the individual endocrine tumour type. Orthopaedic management of fibrous dysplasia may include bisphosphonates and surgical correction of deformities.

\medterm{Medullary Thyroid Carcinoma}-- Neuroendocrine tumor arising from parafollicular C cells that produce calcitonin. Accounts for approximately five percent of thyroid carcinomas. May be sporadic or hereditary as part of multiple endocrine neoplasia type two A or two B. Serum calcitonin serves as a tumor marker. Diagnosis requires FNA with calcitonin immunohistochemistry. Total thyroidectomy with central lymph node dissection is the standard treatment. RET proto-oncogene mutation testing is mandatory in all patients to identify hereditary medullary thyroid carcinoma and guide prophylactic thyroidectomy in affected family members. Tyrosine kinase inhibitors such as vandetanib and cabozantinib are used for advanced or metastatic disease.

\medterm{Menopausal Hormone Therapy (MHT)}-- Menopausal hormone therapy (MHT), formerly called hormone replacement therapy (HRT), is the use of estrogen with or without progestogen to alleviate moderate-to-severe menopausal vasomotor symptoms (hot flashes, night sweats), genitourinary syndrome of menopause (vaginal dryness, dyspareunia, urinary urgency), and prevent osteoporosis. Estrogen-alone therapy is used in post-hysterectomy women; combined therapy (estrogen + progestogen) is required for those with an intact uterus to prevent endometrial hyperplasia/cancer. Regimens: systemic (oral, transdermal patch/gel, spray) for vasomotor symptoms; low-dose vaginal (cream, tablet, ring) for genitourinary symptoms with minimal systemic absorption. The Women's Health Initiative (WHI) established that MHT increases risk of breast cancer, venous thromboembolism, and stroke, though absolute risk is low in women aged <60 or within 10 years of menopause onset. Treatment should use the lowest effective dose for the shortest duration needed, individualizing based on risk profile (cardiovascular, thrombotic, breast cancer).

\medterm{Menopause and Perimenopause}-- Menopause is the permanent cessation of menstruation for 12 consecutive months, reflecting ovarian follicle depletion and estrogen decline, occurring at a median age of 51--52 years in North America. Perimenopause (menopausal transition) begins 4--8 years before the final menstrual period, characterized by variable cycle length, vasomotor symptoms (hot flashes, night sweats—affecting ~80\% of women), sleep disturbance, mood changes, vaginal dryness, and declining fertility. Diagnosis is clinical based on age and menstrual pattern changes; FSH elevation (>25 IU/L) and AMH decline confirm ovarian aging but are not required. Health implications include accelerated bone loss (osteoporosis risk), increased cardiovascular risk, weight gain (central adiposity), and genitourinary syndrome. Management: lifestyle modification (cooling layers, exercise, calcium/vitamin D), MHT for moderate--severe symptoms (most effective), non-hormonal options (SSRIs/SNRIs, gabapentin, oxybutynin) for vasomotor symptoms, and vaginal estrogen for urogenital atrophy.

\medterm{Metabolic Acidosis}-- A primary acid-base disorder characterized by decreased serum bicarbonate (<22 mEq/L) and pH <7.35. Causes are classified by anion gap (AG). High AG: diabetic ketoacidosis, lactic acidosis, renal failure, starvation, toxins (methanol, ethylene glycol, salicylates). Normal AG (hyperchloremic): diarrhea, renal tubular acidosis, acetazolamide. Compensation: respiratory alkalosis (Kussmaul respirations). Treatment: address the underlying cause; sodium bicarbonate reserved for severe acidosis (pH <7.0).

\medterm{Metabolic Syndrome}-- Cluster of interrelated metabolic risk factors increasing cardiovascular disease and type two diabetes risk. Defined by the presence of at least three of five criteria: waist circumference greater than 40 inches in men or greater than 35 inches in women, triglycerides greater than 150 milligrams per deciliter, HDL cholesterol less than 40 milligrams per deciliter in men or less than 50 in women, blood pressure greater than 130/85 millimeters of mercury, and fasting glucose greater than 100 milligrams per deciliter, with treatment thresholds and targets individualised based on overall cardiovascular risk profile and presence of other metabolic risk factors.

\medterm{Metformin in Prediabetes}-- Metformin reduces progression from prediabetes to type two diabetes by approximately thirty-one percent in high-risk individuals. Most effective in younger patients under sixty with BMI greater than 35 and women with prior gestational diabetes. The Diabetes Prevention Program demonstrated sustained benefit over ten years of follow-up. Recommended by many guidelines as an adjunct to lifestyle modification in high-risk prediabetic individuals.

\medterm{Mitochondrial Diabetes}-- Diabetes caused by mutations in mitochondrial DNA particularly the m.3243A>G mutation in the MT-TL1 gene. Maternally inherited. Often associated with sensorineural hearing loss termed maternally inherited diabetes and deafness. Patients are typically lean with impaired insulin secretion. Diagnosis requires genetic testing. Insulin therapy is usually necessary though some patients respond to sulfonylureas. Screening for associated conditions including cardiac disease, seizures, and gastrointestinal dysmotility, and psychiatric disturbance. Genetic counselling is important given the maternal inheritance pattern.

\medterm{Morbid Obesity}-- Severe obesity defined clinically as BMI \(\ge\)40 kg/m\(^2\) or BMI \(\ge\)35 kg/m\(^2\) with obesity-related comorbidities (type 2 diabetes, hypertension, obstructive sleep apnoea, non-alcoholic fatty liver disease, osteoarthritis, cardiovascular disease). Morbid obesity significantly increases all-cause mortality (hazard ratio ~2.5--3.0 compared with normal weight). Management: intensive lifestyle intervention, pharmacotherapy (GLP-1 receptor agonists—semaglutide 2.4 mg weekly, tirzepatide; phentermine--topiramate; naltrexone--bupropion), and bariatric/metabolic surgery (Roux-en-Y gastric bypass, sleeve gastrectomy, adjustable gastric band, biliopancreatic diversion with duodenal switch). Bariatric surgery achieves 25--35\% total body weight loss and induces remission of type 2 diabetes in 60--80\% of patients (the SOS study and STAMPEDE trial). Eligibility per NICE guidelines: BMI \(\ge\)40, or BMI \(\ge\)35 with recent-onset type 2 diabetes; consideration of psychological suitability, commitment to lifelong follow-up, and nutritional supplementation.

\medterm{Multiple Endocrine Neoplasia Type 1}-- An autosomal dominant tumor syndrome caused by MEN1 gene mutation (menin). Characterized by tumors of the parathyroid (hyperparathyroidism in 95% of patients by age 50, most common), pancreatic islet cells (gastrinoma causing Zollinger-Ellison syndrome, insulinoma, glucagonoma, VIPoma), and anterior pituitary (prolactinoma most common, GH, ACTH). Other tumors: carcinoid, adrenal adenoma, lipomas. Diagnosis: clinical criteria, genetic testing for MEN1 mutation. Management: surgical resection of tumors; peptide receptor radionuclide therapy for advanced neuroendocrine tumors.

\medterm{Multiple Endocrine Neoplasia Type 2}-- An autosomal dominant tumor syndrome caused by RET proto-oncogene mutations. Type 2A (95%): medullary thyroid carcinoma (MTC, nearly 100% penetrance), pheochromocytoma (50%), and hyperparathyroidism (20-30%). Type 2B: MTC, pheochromocytoma, mucosal neuromas, marfanoid habitus; more aggressive MTC. Management: prophylactic thyroidectomy based on RET mutation risk stratification; screen for pheochromocytoma before thyroid surgery; annual lab and imaging surveillance.

\medterm{Myxedema Coma}-- Severe decompensated hypothyroidism characterized by altered mental status, hypothermia, bradycardia, hypoventilation, hyponatremia, and hypoglycemia. Most commonly occurs in elderly patients with long-standing untreated hypothyroidism precipitated by infection, surgery, or sedative medications. Mortality rates exceed fifty percent. Diagnosis is clinical with low free T4 and elevated TSH. Treatment includes IV levothyroxine and IV hydrocortisone to prevent adrenal crisis, passive rewarming, and supportive care in the intensive care unit setting including respiratory and haemodynamic support.

\medterm{Neonatal Diabetes Mellitus}-- Diabetes diagnosed within the first six months of life. Can be transient or permanent. Transient neonatal diabetes mellitus is caused by abnormalities at chromosome 6q24 and usually resolves but may relapse in adolescence. Permanent neonatal diabetes mellitus is most commonly caused by mutations in the KCNJ11 or ABCC8 genes encoding the KATP channel subunits. Patients with KCNJ11 mutations can often transition from insulin to sulfonylureas with excellent glycemic control due to the specific mechanism of action targeting the KATP channel defect.

\medterm{Neuroendocrine Tumors Grading}-- Classification of neuroendocrine tumors by proliferative rate using Ki-67 index and mitotic count. Well-differentiated low-grade tumors have Ki-67 less than three percent and fewer than two mitoses per ten high-power fields. Intermediate-grade tumors have Ki-67 of three to twenty percent. Poorly differentiated high-grade tumors have Ki-67 greater than twenty percent. Grading guides prognosis and treatment intensity with low-grade tumors treated with somatostatin analogs and high-grade tumors require chemotherapy including everolimus and temozolomide and peptide receptor radionuclide therapy.

\medterm{Neurofibromatosis Type 1}-- Autosomal dominant disorder caused by mutations in the NF1 gene encoding neurofibromin on chromosome seventeen. Diagnostic criteria include two or more of the following: six or more cafe-au-lait spots, two or more neurofibromas or one plexiform neurofibroma, axillary or inguinal freckling, optic pathway glioma, two or more Lisch nodules, bony abnormalities, and a first-degree relative with NF1. Endocrine manifestations include pheochromocytoma and growth hormone excess.

\medterm{Non-Functioning Pituitary Adenoma}-- Pituitary adenomas that do not secrete clinically significant hormones. The most common pituitary adenomas to present with mass effects. Typically macroadenomas at diagnosis causing headaches, visual field defects particularly bitemporal hemianopia from optic chiasm compression, and hypopituitarism from compression of normal pituitary tissue. Diagnosis is by MRI showing a pituitary mass and hormone testing to exclude functioning adenomas. Management includes transsphenoidal surgery for symptomatic macroadenomas causing visual compromise or hypopituitarism. Post-operative surveillance is required with MRI and hormone monitoring to detect recurrence.

\medterm{Normocalcemic Hyperparathyroidism}-- Condition of elevated parathyroid hormone with persistently normal serum calcium and normal vitamin D levels. Exclusion of secondary causes of elevated PTH including vitamin D deficiency, renal insufficiency, and medications is required. Represents an early or mild form of primary hyperparathyroidism. May progress to overt hypercalcemia. Monitoring includes annual calcium PTH levels and bone density assessment. Parathyroidectomy is considered for patients with complications or progressive disease. Management includes monitoring calcium, PTH, and bone density annually, with parathyroidectomy considered for patients who develop complications such as worsening bone density, nephrolithiasis, or progression to hypercalcaemia.

\medterm{Obesity}-- Abnormal or excessive fat accumulation that presents a health risk. Defined by BMI: \(\ge\)30 kg/m² (class I: 30-34.9, class II: 35-39.9, class III: \(\ge\)40). Major risk factor for type 2 diabetes, hypertension, cardiovascular disease, sleep apnea, osteoarthritis, certain cancers, and reduced life expectancy. Management: lifestyle interventions, pharmacotherapy (GLP-1 agonists, phentermine-topiramate, naltrexone-bupropion), and bariatric surgery for severe obesity.

\medterm{Oestrogen Replacement Therapy}-- See Hormone Replacement Therapy (HRT).

\medterm{Osteomalacia}-- Impaired mineralization of newly formed osteoid in adults resulting from vitamin D deficiency, phosphate depletion, or defective mineralization. Causes include vitamin D deficiency from inadequate intake, malabsorption, and impaired synthesis, phosphate depletion from tumor-induced osteomalacia with fibroblast growth factor twenty-three excess, and medications including bisphosphonates and aluminum. Presents with diffuse bone pain, muscle weakness particularly proximal, and increased fracture risk. Diagnosis is confirmed by low 25-hydroxyvitamin D levels, elevated alkaline phosphatase, and characteristic radiographic findings including Looser zones (pseudofractures) on X-ray. Treatment is high-dose vitamin D supplementation with calcium, and treatment of the underlying cause where possible.

\medterm{Papillary Thyroid Carcinoma}-- Most common histological subtype of thyroid carcinoma accounting for approximately eighty percent of cases. Characterized by nuclear features including nuclear grooves, intranuclear cytoplasmic inclusions, and Orphan Annie eye nuclei. Often multifocal and bilateral with lymph node metastases. Excellent prognosis with greater than ninety-five percent ten-year survival. Treatment includes total thyroidectomy with radioactive iodine ablation for intermediate and high-risk disease. TSH suppression with levothyroxine for high-risk patients to reduce recurrence risk.

\medterm{Parathyroid Crisis}-- Severe hypercalcemia with calcium greater than fourteen milligrams per deciliter causing altered mental status, cardiac arrhythmias, and acute kidney injury. Most commonly occurs in patients with primary hyperparathyroidism precipitated by dehydration, immobility, or thiazide diuretics. Treatment requires aggressive IV normal saline hydration, loop diuretics after adequate hydration, calcitonin for rapid but transient calcium reduction, bisphosphonates for sustained calcium lowering, and denosumab for refractory cases. Urgent parathyroidectomy is the definitive treatment once the patient is stabilised.

\medterm{Parathyroid Hormone Physiology}-- Parathyroid hormone is the primary regulator of calcium homeostasis secreted by the four parathyroid glands in response to low ionized calcium levels detected by the calcium-sensing receptor. PTH increases serum calcium through three mechanisms: stimulating osteoclastic bone resorption to release calcium and phosphorus, increasing renal calcium reabsorption in the distal tubule, and stimulating renal production of active vitamin D which enhances intestinal calcium absorption. PTH also increases phosphate excretion in the proximal tubule, thereby lowering serum phosphate. PTH secretion is tightly regulated by ionised calcium through negative feedback on the calcium-sensing receptor.

\medterm{PCOS}-- See Polycystic Ovary Syndrome.

\medterm{Pheochromocytoma}-- Catecholamine-secreting tumor arising from chromaffin cells of the adrenal medulla or extra-adrenal paraganglia. Approximately ten percent are bilateral, ten percent are extra-adrenal, and ten percent are malignant. Classic presentation includes episodic headache, sweating, and palpitations with sustained or paroxysmal hypertension. Associated with hereditary syndromes including von Hippel-Lindau, multiple endocrine neoplasia type two, and neurofibromatosis type one. Diagnosis is by plasma free metanephrines or 24-hour urinary fractionated metanephrines, which have high sensitivity for detection. Imaging with CT or MRI of the abdomen and pelvis localises the tumour. MIBG scintigraphy or PET-CT with FDG or gallium-DOTATATE may be used for staging and detection of metastatic disease. Definitive treatment is laparoscopic adrenalectomy following adequate preoperative alpha-blockade with phenoxybenzamine or doxazosin for at least ten to fourteen days, combined with beta-blockade after alpha-blockade is established. Metastatic disease requires debulking surgery, radionuclide therapy, or tyrosine kinase inhibitors.

\medterm{Phosphate Homeostasis}-- Serum phosphate is regulated by the interplay of PTH vitamin D and fibroblast growth factor 23. PTH promotes renal phosphate excretion. Active vitamin D stimulates intestinal phosphate absorption. FGF23 produced by osteocytes inhibits renal phosphate reabsorption and suppresses 1-alpha-hydroxylase reducing active vitamin D production. Hyperphosphatemia occurs in advanced CKD and is associated with cardiovascular calcification and increased mortality. Management includes dietary phosphate restriction and phosphate binders in chronic kidney disease.

\medterm{Pituitary Apoplexy}-- Acute hemorrhage or infarction within a pituitary adenoma causing sudden headache, visual field defects, ophthalmoplegia, and altered consciousness. Medical emergency requiring urgent glucocorticoid administration to prevent adrenal crisis from acute ACTH deficiency. Severe cases with visual compromise require emergent transsphenoidal surgery. Mild cases may be managed conservatively with close monitoring and hormone replacement. Risk factors include pregnancy, dopamine agonist therapy, and anticoagulation as risk factors. Immediate management focuses on haemodynamic stabilisation, glucocorticoid administration, and urgent neurosurgical consultation.

\medterm{Polycystic Ovary Syndrome}-- Common endocrine disorder affecting six to twelve percent of reproductive-aged women characterized by hyperandrogenism, ovulatory dysfunction, and polycystic ovarian morphology on ultrasound. Diagnostic criteria include two of three features: oligo-anovulation, clinical or biochemical hyperandrogenism, and polycystic ovaries after exclusion of other causes. Associated with insulin resistance, obesity, type two diabetes, and cardiovascular risk. Treatment includes lifestyle modification for weight management, combined oral contraceptive or progestin therapy for menstrual regulation and endometrial protection, antiandrogens like spironolactone for hirsutism, metformin for insulin resistance and glucose intolerance, and ovulation induction with letrozole or clomiphene citrate for infertility. Long-term monitoring for cardiovascular risk factors and glucose intolerance is essential.

\medterm{Precocious Puberty}-- Onset of secondary sexual characteristics before age eight in girls and age nine in boys. Central precocious puberty results from early activation of the hypothalamic-pituitary-gonadal axis. Causes include CNS tumors, hydrocephalus, and idiopathic. Peripheral precocious puberty results from autonomous sex steroid production from adrenal tumors, gonadal tumors, or congenital adrenal hyperplasia. Diagnosis includes bone age X-ray, GnRH stimulation test, and imaging of the brain and adrenal glands. Treatment is with GnRH agonists for central precocious puberty to suppress the axis and slow pubertal progression, while peripheral causes are managed by addressing the underlying source of sex steroid production.

\medterm{Prediabetes}-- Intermediate state of hyperglycemia between normal glucose homeostasis and diabetes defined by fasting glucose 100 to 125 milligrams per deciliter, two-hour glucose during oral glucose tolerance test 140 to 199 milligrams per deciliter, or A1c 5.7 to 6.4 percent. Individuals are at increased risk for progression to type two diabetes with annual progression rates of approximately five to ten percent. Cardiovascular disease risk is already elevated. Lifestyle intervention with weight loss of five to seven percent and increased physical activity is the cornerstone of prevention and is more effective than metformin. For individuals with additional risk factors, metformin may be considered.

\medterm{Premature Ovarian Insufficiency}-- Loss of ovarian function before age forty characterized by amenorrhea, hypergonadotropic hypogonadism, and estrogen deficiency. Causes include genetic disorders particularly Turner syndrome and fragile X premutations, autoimmune oophoritis, iatrogenic from chemotherapy or radiation, and idiopathic. Associated with decreased bone density, cardiovascular risk, and infertility. Diagnosis requires four months of amenorrhea with elevated FSH greater than twenty-five milli-international units per millilitre on two occasions four weeks apart. Treatment includes hormone replacement therapy with estrogen and progesterone, and fertility options including oocyte donation for those desiring pregnancy.

\medterm{Pretibial Myxedema}-- Localized dermopathy occurring in approximately five percent of patients with Graves disease caused by glycosaminoglycan deposition in the dermis particularly over the pretibial area. Presents as non-pitting indurated plaques with a peau dorange appearance. Pathogenesis involves orbital fibroblasts stimulated by thyroid autoantibodies producing excessive hyaluronic acid. Usually occurs in patients with severe ophthalmopathy. Treatment includes topical corticosteroids and compression bandaging through a multidisciplinary approach with dermatology and endocrinology input.

\medterm{Primary Aldosteronism}-- Overproduction of aldosterone by the adrenal glands (Conn syndrome: adrenal adenoma; or bilateral adrenal hyperplasia). Causes hypertension with hypokalemia (40% normokalemic). Presents with treatment-resistant hypertension with or without hypokalemia, metabolic alkalosis, low plasma renin, and high aldosterone-to-renin ratio (ARR >30). Diagnosis: elevated ARR, confirmatory testing (saline suppression test, captopril challenge, oral sodium loading), adrenal CT and adrenal vein sampling to distinguish unilateral from bilateral. Treatment: unilateral adrenalectomy for adenoma; mineralocorticoid receptor antagonists (spironolactone, eplerenone) for bilateral hyperplasia.

\medterm{Primary Hyperaldosteronism}-- Autonomous aldosterone secretion most commonly caused by an aldosterone-producing adrenal adenoma termed Conn syndrome. Leads to hypertension, hypokalemia, metabolic alkalosis, and suppressed renin. Accounts for approximately five to ten percent of hypertensive patients. Screening is by elevated aldosterone-to-renin ratio. Diagnosis requires confirmatory salt loading test demonstrating failure to suppress aldosterone. Localization is by adrenal CT and adrenal vein sampling to distinguish unilateral from bilateral disease. Unilateral aldosterone-producing adenoma is treated with laparoscopic adrenalectomy. Bilateral hyperplasia is managed medically with mineralocorticoid receptor antagonists such as spironolactone or eplerenone.

\medterm{Primary Hyperparathyroidism}-- Autonomous overproduction of parathyroid hormone from one or more abnormal parathyroid glands causing hypercalcemia. Most commonly caused by a single parathyroid adenoma accounting for approximately eighty-five percent. Multigland hyperplasia accounts for approximately ten to fifteen percent and parathyroid carcinoma is rare. Clinical manifestations include nephrolithiasis, osteoporosis, neuropsychiatric symptoms, gastrointestinal symptoms, and cardiovascular disease. Many patients are asymptomatic at diagnosis, identified incidentally through routine biochemical testing. Surgery is recommended for all symptomatic patients and for asymptomatic patients who meet criteria including age under fifty, serum calcium greater than one milligram per decilitre above the upper limit of normal, osteoporosis, nephrolithiasis, or renal impairment.

\medterm{Prolactin Physiology}-- Prolactin is the only pituitary hormone primarily under tonic inhibitory control by the hypothalamus through dopamine acting on D2 receptors. Prolactin stimulates lactation suppresses GnRH and is important for reproductive function. Physiological elevations occur during pregnancy lactation breast stimulation sleep and stress. Hyperprolactinemia suppresses the hypothalamic-pituitary-gonadal axis causing amenorrhea infertility and galactorrhea in women and erectile dysfunction and hypogonadism in men. Management of hyperprolactinemia involves treatment of the underlying cause and dopamine agonist therapy when indicated.

\medterm{Prolactinoma Treatment}-- Dopamine agonists are first-line medical therapy for prolactinomas. Cabergoline is preferred over bromocriptine due to higher efficacy, less side effects, and twice-weekly dosing improving compliance. Normalize prolactin levels in approximately eighty to ninety percent of patients and significantly reduce tumor size in most. Surgery is reserved for dopamine agonist-resistant tumors, patient intolerance, or cerebrospinal fluid leak from macroadenomas. Radiation therapy is reserved for refractory cases. Long-term monitoring with serial prolactin measurements and MRI is required to detect recurrence.

\medterm{Radioactive Iodine Ablation}-- Treatment of differentiated thyroid carcinoma using iodine-131 to destroy residual thyroid tissue after total thyroidectomy. Indications include intermediate and high-risk disease for remnant ablation and treatment of residual or recurrent disease. Patients must be hypothyroid or receive recombinant TSH to maximize iodine uptake. Preparation includes low-iodine diet for one to two weeks and discontinuation of thyroid hormone or levothyroxine withdrawal. Side effects include sialadenitis, taste changes, xerostomia, and rarely radiation thyroiditis. Long-term risks include secondary malignancies such as leukaemia and salivary gland tumours. Gonadal dysfunction may occur, and sperm or oocyte cryopreservation should be considered before treatment. Pregnancy is contraindicated for six to twelve months after ablation.

\medterm{Renin-Angiotensin-Aldosterone System Physiology}-- The RAAS regulates blood pressure and fluid electrolyte balance. Renin released from juxtaglomerular cells cleaves angiotensinogen to angiotensin I. Angiotensin converting enzyme converts angiotensin I to angiotensin II which causes vasoconstriction and stimulates aldosterone secretion from the adrenal cortex. Aldosterone acts on the collecting duct to promote sodium reabsorption and potassium secretion. The system is activated by decreased renal perfusion decreased sodium delivery to the macula densa, and decreased sodium chloride delivery. Dysregulation of this system contributes to hypertension, heart failure, and chronic kidney disease. Pharmacological inhibitors of the RAAS including ACE inhibitors, angiotensin receptor blockers, and mineralocorticoid receptor antagonists are cornerstones of cardiovascular and renal protective therapy.

\medterm{Rickets}-- A childhood disorder of impaired mineralization of growing bone due to vitamin D deficiency, calcium deficiency, or phosphate deficiency. Presents with bone pain, deformities (bowed legs, wrist enlargement, rachitic rosary), growth retardation, and hypocalcemic seizures. Predisposing factors: dark skin, limited sun exposure, exclusive breastfeeding without supplementation. Diagnosis: low 25-hydroxyvitamin D, elevated ALP, characteristic X-ray findings. Treatment: high-dose vitamin D and calcium supplementation.

\medterm{Sheehan Syndrome}-- Postpartum pituitary necrosis resulting from severe hemorrhage and hypotension during delivery. The pituitary gland is enlarged during pregnancy and particularly vulnerable to ischemic injury from blood loss. Presentation varies from fatigue and failure to lactate to panhypopituitarism. Growth hormone and prolactin deficiency occur first. Diagnosis is by hormone testing. Treatment is lifelong hormone replacement of all deficient axes.

\medterm{Statins in Diabetes}-- HMG-CoA reductase inhibitors are first-line therapy for cardiovascular risk reduction in diabetic patients. Most adults with diabetes aged forty to seventy-five benefit from moderate-intensity statin therapy regardless of baseline LDL cholesterol. High-intensity statin therapy with atorvastatin forty to eighty milligrams or rosuvastatin twenty to forty milligrams is indicated for patients with established atherosclerotic cardiovascular disease or multiple risk factors. Benefits include reduction in major adverse cardiovascular events and all-cause mortality. Statin therapy should be continued regardless of baseline lipid levels, and adherence is critical for optimal cardiovascular prevention.

\medterm{Subacute Thyroiditis}-- Inflammatory condition of the thyroid most likely viral in etiology presenting with painful tender thyroid enlargement, fever, and malaise. Associated with elevated erythrocyte sedimentation rate and C-reactive protein. Disease course includes thyrotoxicosis from hormone release during gland destruction, followed by hypothyroidism, and then recovery in most patients. Diagnosis is clinical with the typical presentation and characteristic elevated inflammatory markers. Treatment includes NSAIDs for pain and inflammation, and beta-blockers for symptomatic thyrotoxicosis during the initial phase. Glucocorticoids may be used in severe cases. Thyroid function should be monitored as the patient may require levothyroxine replacement during the subsequent hypothyroid phase.

\medterm{Teriparatide for Osteoporosis}-- Recombinant human parathyroid hormone fragments one through thirty-four that stimulate osteoblast-mediated bone formation. Indicated for severe osteoporosis with fractures or very low T-scores and for patients who have failed or are intolerant to bisphosphonates. Daily subcutaneous injection for up to two years. Increases bone density more rapidly than anti-resorptive agents. Must be followed by anti-resorptive therapy with bisphosphonates or denosumab to maintain gains. Romosozumab an anti-sclerostin antibody, builds bone rapidly and is given as monthly subcutaneous injections for one year, after which antiresorptive therapy is required to maintain bone density gains.

\medterm{Testosterone}-- Testosterone is the primary male sex hormone, an androgen produced mainly in the testicles that is essential for the development of male reproductive tissues, secondary sexual characteristics, and overall health. In addition to its role in sexual function and sperm production, testosterone influences muscle mass, bone density, fat distribution, red blood cell production, and mood. Testosterone levels naturally decline with age, and abnormally low levels can cause reduced libido, erectile dysfunction, fatigue, loss of muscle mass, and osteoporosis; testosterone replacement therapy is available but requires careful consideration of potential risks.

\medterm{Tetany}-- A clinical syndrome of neuromuscular hyperexcitability caused by hypocalcaemia (most common—total calcium <2.1 mmol/L, ionised calcium <1.1 mmol/L), but also caused by hypomagnesaemia, hypokalaemia, and alkalosis (respiratory or metabolic). Symptoms: perioral paraesthesias (tingling around the mouth), numbness and tingling in the fingers and toes, muscle cramps, carpopedal spasm (flexion of the wrist and metacarpophalangeal joints, extension of the interphalangeal joints, adduction of the thumb—'obstetrician's hand'), and laryngospasm (stridor, respiratory distress—life-threatening). Latent tetany: elicited by Trousseau sign (carpopedal spasm induced by inflation of a BP cuff above systolic pressure for 3 minutes) and Chvostek sign (tapping the facial nerve anterior to the ear causes ipsilateral facial muscle twitching). Causes of hypocalcaemia: hypoparathyroidism (post-surgical—most common, autoimmune, DiGeorge syndrome), vitamin D deficiency/insufficiency, chronic kidney disease, acute pancreatitis, pseudohypoparathyroidism, and drugs (bisphosphonates, denosumab, furosemide, phenytoin). Treatment: IV calcium gluconate 10--20 mL of 10% solution (for acute symptomatic tetany) given slowly with cardiac monitoring, followed by oral calcium and vitamin D/calcitriol. Treat underlying cause.

\medterm{Thiazolidinediones}-- Insulin sensitizers activating peroxisome proliferator-activated receptors primarily in adipose tissue and skeletal muscle. Agents include pioglitazone and rosiglitazone. Mechanisms include enhanced insulin-mediated glucose uptake, altered adipokine secretion, and reduced hepatic glucose output. Effective A1c reduction of one to one point five percent. Side effects include weight gain, fluid retention leading to heart failure exacerbation, increased fracture risk particularly in postmenopausal women and in premenopausal women with polycystic ovary syndrome where they may promote ovulation. Due to safety concerns, pioglitazone is more commonly used than rosiglitazone, and both require careful monitoring for adverse effects.

\medterm{Thyroglobulin Surveillance}-- Thyroglobulin produced only by thyroid follicular cells serves as a tumor marker after total thyroidectomy for differentiated thyroid carcinoma. In patients who have undergone complete thyroidectomy and radioactive iodine ablation thyroglobulin should be undetectable. Rising thyroglobulin levels indicate residual or recurrent thyroid tissue. Stimulated thyroglobulin after recombinant human TSH or thyroid hormone withdrawal provides greater sensitivity for detecting recurrence. The thyroglobulin antibody assay is essential as autoantibodies can interfere with the measurement causing falsely low or undetectable levels, limiting the utility of the test in antibody-positive patients.

\medterm{Thyroglossal Duct Cyst}-- Congenital midline neck mass resulting from persistence of the thyroglossal duct the embryological tract of thyroid descent. Most common congenital neck mass presenting in childhood. Located at or near the hyoid bone and moves with swallowing and tongue protrusion. Diagnosis is by ultrasound demonstrating a cystic midline neck mass. Thyroid tissue may be present within the cyst. Treatment is surgical excision with Sistrunk procedure which includes removal of the central portion of the hyoid bone to prevent recurrence related to the embryological tract. Recurrence is rare after complete excision.

\medterm{Thyroid Cancer Staging}-- American Joint Committee on Cancer TNM staging for differentiated thyroid carcinoma incorporates tumor size, extent of invasion, lymph node metastases, and distant metastases. Age is a critical prognostic factor with patients under fifty having more favorable outcomes. Stage I is any tumor without metastases in patients under fifty. Stage II includes distant metastases or extrathyroidal extension in patients under fifty and larger tumors in patients over fifty. Risk stratification systems include the American Thyroid Association risk stratification, which categorises patients into low, intermediate, and high-risk groups to guide post-operative management including the need for radioactive iodine ablation and the intensity of TSH suppression therapy.

\medterm{Thyroid Disorders}-- Thyroid disorders are conditions affecting the thyroid gland, a butterfly-shaped endocrine organ in the neck that produces hormones regulating metabolism, growth, and development. Hyperthyroidism involves excessive thyroid hormone production, causing symptoms such as weight loss, tachycardia, heat intolerance, and anxiety, while hypothyroidism results from insufficient hormone production and leads to fatigue, weight gain, cold intolerance, and depression. Autoimmune diseases including Graves disease and Hashimoto thyroiditis are common causes; management includes medications, radioactive iodine therapy, or thyroid surgery depending on the specific disorder.

\medterm{Thyroid Eye Disease}-- Autoimmune inflammatory disorder of the orbital tissues associated with Graves disease characterized by expansion of extraocular muscles and orbital fat. Pathogenesis involves orbital fibroblasts expressing the TSH receptor targeted by the same thyroid-stimulating immunoglobulins affecting the thyroid. Clinical features include proptosis, periorbital edema, chemosis, lid retraction, and diplopia. Severe disease can cause optic nerve compression. Treatment includes smoking cessation, selenium supplementation which has been shown to slow progression of mild to moderate thyroid eye disease, glucocorticoids for active inflammatory disease, and orbital radiotherapy or surgical decompression for sight-threatening proptosis or optic neuropathy. Teprotumumab, an IGF-1 receptor inhibitor, is a newer targeted therapy that reduces proptosis and diplopia in active disease.

\medterm{Thyroid Function Tests}-- Laboratory assessment of the hypothalamic-pituitary-thyroid axis. TSH is the most sensitive screening test with elevated levels indicating primary hypothyroidism and suppressed levels indicating hyperthyroidism. Free T4 measures the biologically active unbound thyroxine. Total T4 and T3 reflect bound and unbound hormone and are affected by binding protein abnormalities. T3 thyrotoxicosis occurs when T3 is preferentially elevated in early Graves disease or toxic adenoma. TSH receptor antibodies can also be measured to confirm autoimmune aetiology. Interpretation of thyroid function tests must account for non-thyroidal illness, medication effects, and binding protein abnormalities that can alter total hormone levels without affecting free hormone concentrations.

\medterm{Thyroid Nodule Evaluation}-- Systematic approach to evaluation of thyroid nodules detected by physical examination or imaging. Thyroid ultrasound characterizes nodule features including size, composition, echogenicity, margins, calcifications, and shape. TI-RADS classification stratifies malignancy risk and guides biopsy decisions. Fine-needle aspiration biopsy is indicated for nodules meeting size and suspicious feature criteria. Bethesda cytology classification categorizes results from non-diagnostic through suspicious for malignancy to malignant. Management of benign nodules is observation with serial ultrasound, while malignant or suspicious nodules require surgical resection. Molecular testing of biopsy specimens may help refine risk stratification and guide surgical decision-making.

\medterm{Thyroid Storm}-- Life-threatening exacerbation of thyrotoxicosis presenting with hyperpyrexia, tachycardia, heart failure, altered mental status, and multi-organ dysfunction. Precipitated by infection, surgery, trauma, or iodine load in patients with untreated hyperthyroidism. Mortality rates range from twenty to thirty percent. Diagnosis is clinical based on Burch-Wartofsky scoring system. Treatment includes beta-blockers to control heart rate, thionamides to block new thyroid hormone synthesis, iodine to block thyroid hormone release (the Plummer effect), corticosteroids to reduce peripheral conversion of T4 to T3 and support haemodynamics, and aggressive supportive care including cooling, hydration, and management of heart failure. Definitive treatment of hyperthyroidism with radioactive iodine or thyroidectomy should follow after stabilisation.

\medterm{Thyroiditis}-- Inflammation of the thyroid gland with various etiologies. Hashimoto thyroiditis: autoimmune, most common cause of hypothyroidism in iodine-sufficient areas. Subacute (De Quervain) thyroiditis: post-viral, presents with painful thyroid and transient hyperthyroidism followed by hypothyroidism. Postpartum thyroiditis: autoimmune, occurs within 1 year of delivery. Silent thyroiditis: painless lymphocytic thyroiditis. Diagnosis: TSH, T4, TPO antibodies, thyroid ultrasound, and radioactive iodine uptake.

\medterm{Thyroiditis Classification}-- Inflammatory conditions of the thyroid classified by clinical presentation and etiology. Hashimoto thyroiditis is chronic autoimmune with hypothyroidism. Graves disease is autoimmune with hyperthyroidism. Subacute thyroiditis is painful and likely viral. Silent thyroiditis and postpartum thyroiditis are painless autoimmune causing transient thyrotoxicosis followed by hypothyroidism. Riedel thyroiditis is rare idiopathic fibrosis causing a stony-hard goiter that may compress adjacent structures. Diagnosis is by clinical presentation, thyroid function tests, ultrasound, and radioactive iodine uptake which is low in inflammatory thyroiditis and high in Graves disease. Management is directed at the underlying aetiology.

\medterm{Time in Range}-- Percentage of time that continuous glucose monitoring readings remain within the target glucose range of seventy to one hundred eighty milligrams per deciliter for most adults with diabetes. General target is greater than seventy percent corresponding to approximately seventeen hours per day. Time below range less than seventy is less than four percent. Time below fifty-four is less than one percent. Time above range greater than one hundred eighty is less than twenty-five percent. Time in range correlates well with A1c and is a more dynamic measure of glycaemic control. Each ten percent increase in time in range corresponds to an approximate 0.5 percent reduction in A1c. Time in range targets may be individualised for older adults, pregnant women, and those with frequent hypoglycaemia.

\medterm{Turner Syndrome}-- Chromosomal abnormality in females with complete or partial absence of one X chromosome most commonly with a karyotype of 45,X occurring in one in two thousand five hundred live female births. Features include short stature, webbed neck, shield chest, lymphedema, streak gonads causing primary amenorrhea and infertility, congenital heart defects particularly bicuspid aortic valve and coarctation of the aorta, and renal anomalies. Ovarian failure leads to estrogen deficiency requiring hormone replacement therapy for breast development and maintenance of secondary sexual characteristics, and oestrogen replacement for bone health and cardiovascular protection. Growth hormone therapy may be considered during childhood to improve adult height. Fertility options including oocyte donation should be discussed.

\medterm{Type 1 Diabetes Mellitus}-- Type 1 diabetes is an autoimmune disease in which the immune system attacks and destroys the insulin-producing beta cells of the pancreas, resulting in absolute insulin deficiency. It typically presents in childhood or adolescence with symptoms of polyuria, polydipsia, weight loss, and fatigue, and requires lifelong insulin therapy through multiple daily injections or an insulin pump. Management involves careful monitoring of blood glucose levels, carbohydrate counting, and balancing insulin doses with physical activity to maintain near-normal glycemia and prevent acute complications like diabetic ketoacidosis and long-term complications affecting the eyes, kidneys, nerves, and cardiovascular system.

\medterm{Type 2 Diabetes Mellitus}-- Type 2 diabetes is a chronic metabolic disorder characterized by insulin resistance and progressive pancreatic beta-cell dysfunction, leading to hyperglycemia. It is strongly associated with obesity, physical inactivity, and genetic predisposition, and typically develops in adulthood, though increasing rates are seen in younger populations. Management begins with lifestyle modifications including weight loss, dietary changes, and exercise, then progresses to oral medications such as metformin, and may eventually require insulin as beta-cell function declines over time.

\medterm{Vasopressin (Antidiuretic Hormone, ADH)}-- See Antidiuretic Hormone (ADH).

\medterm{VIPoma}-- A rare pancreatic neuroendocrine tumor secreting vasoactive intestinal peptide (VIP), causing WDHA syndrome: watery diarrhea (large volume, 3-5 L/day), hypokalemia, and achlorhydria/acidosis. Most are malignant and located in the pancreatic tail. Diagnosis: elevated VIP levels, pancreatic tumor on imaging. Treatment: surgical resection for localized disease; somatostatin analogs (octreotide) control symptoms; chemotherapy or peptide receptor radionuclide therapy for metastatic disease.

\medterm{Virilisation}-- The development of male secondary sexual characteristics in females, resulting from excessive androgen production or exposure. Causes: congenital adrenal hyperplasia (21-hydroxylase deficiency, the most common), androgen-secreting tumours (ovarian or adrenal), Cushing syndrome, exogenous androgen use (anabolic steroids, testosterone therapy), and polycystic ovary syndrome (mild hyperandrogenism). Clinical features: hirsutism (male-pattern hair growth on face, chest, back, and abdomen), temporal hair recession (male-pattern baldness), deepening of the voice, clitoromegaly, increased muscle mass, acne, and amenorrhoea/oligomenorrhoea. The severity and age of onset depend on the cause: prenatal virilisation (ambiguous genitalia at birth, classically in congenital adrenal hyperplasia) vs. postpubertal virilisation (gradual onset in PCOS, rapid onset suggests tumour). Evaluation: serum testosterone (total and free), DHEAS, 17-hydroxyprogesterone, and imaging (ovarian/adrenal ultrasound, CT/MRI). Treatment: directed at the underlying cause — glucocorticoid replacement for CAH, surgical resection of tumours, antiandrogens (spironolactone, cyproterone acetate) and cosmetic measures (electrolysis, laser hair removal) for PCOS.

\medterm{Von Hippel-Lindau Disease}-- Autosomal dominant syndrome caused by mutations in the VHL tumor suppressor gene on chromosome three predisposing to hemangioblastomas of the central nervous system and retina, clear cell renal cell carcinoma, pheochromocytoma, pancreatic cysts and neuroendocrine tumors, and endolymphatic sac tumors. Retinal hemangioblastomas are the most common initial manifestation. CNS hemangioblastomas occur most commonly in the cerebellum and spinal cord. Screening includes annual ophthalmological examination, imaging of the brain and spine every two to three years, and abdominal imaging for renal cell carcinoma and pancreatic lesions. Treatment is surgical resection of symptomatic tumours, with targeted therapies such as sunitinib and belzutifan for advanced renal cell carcinoma.

\medterm{Wolfram Syndrome}-- Autosomal recessive disorder caused by mutations in the WFS1 gene encoding wolframin. Classic triad of diabetes mellitus, optic atrophy, and diabetes insipidus. Additional features include sensorineural deafness, neurodegeneration, and psychiatric illness. Onset of diabetes typically in the first decade. Progressive optic atrophy leads to blindness. Diagnosis is by clinical criteria and genetic testing. Treatment is supportive with insulin for diabetes and desmopressin for diabetes insipidus.
